
\documentclass{article}
\usepackage{colortbl}
\usepackage{makecell}
\usepackage{multirow}
\usepackage{supertabular}

\begin{document}

\newcounter{utterance}

\centering \large Interaction Transcript for game `clean\_up', experiment `2\_hard\_5obj\_de', episode 3 with Qwen3.6{-}35B{-}A3B{-}FP8{-}without{-}reasoning.
\vspace{24pt}

{ \footnotesize  \setcounter{utterance}{1}
\setlength{\tabcolsep}{0pt}
\begin{supertabular}{c@{$\;$}|p{.15\linewidth}@{}p{.15\linewidth}p{.15\linewidth}p{.15\linewidth}p{.15\linewidth}p{.15\linewidth}}
    \# & $\;$A & \multicolumn{4}{c}{Game Master} & $\;\:$B\\
    \hline

    \theutterance \stepcounter{utterance}  
    & & \multicolumn{4}{p{0.6\linewidth}}{
        \cellcolor[rgb]{0.9,0.9,0.9}{
            \makecell[{{p{\linewidth}}}]{
                \texttt{\tiny{[P1$\langle$GM]}}
                \texttt{Ich bin der Spielleiter, und du spielst ein kooperatives Spiel mit dem folgenden Raster als Spielbrett:} \\
\\ 
\texttt{```} \\
\texttt{ 1234567} \\
\texttt{╔═╤═╤═══╗} \\
\texttt{║C│◌│◌TI║ 1} \\
\texttt{║◌├─┤W◌◌║ 2} \\
\texttt{╟─┤◌├─┬─╢ 3} \\
\texttt{║◌│◌│◌│◌║ 4} \\
\texttt{╟─┼─┘◌│◌║ 5} \\
\texttt{║◌│◌H┌┴─╢ 6} \\
\texttt{║◌│◌◌│◌◌║ 7} \\
\texttt{╚═╧══╧══╝} \\
\\ 
\texttt{```} \\
\\ 
\texttt{* Am oberen Rand sind die x{-}Koordinaten angegeben, die nach rechts hin zunehmen, und am rechten Rand die y{-}Koordinaten, die nach unten zunehmen.} \\
\texttt{* Die folgenden Objekte sind zufällig auf deinem Raster platziert: 'W', 'I', 'T', 'C', 'H'.} \\
\\ 
\\ 
\texttt{Dein Mitspieler sieht eine Variation des Spielbretts, auf dem dieselben Objekte an anderen zufälligen Positionen platziert sind. Du kannst das Spielbrett des anderen Spielers nicht sehen, und er kann deins nicht sehen.} \\
\\ 
\\ 
\texttt{**Ziel des Spiels:**} \\
\\ 
\texttt{Beide Spieler müssen die Objekte auf ihrem jeweiligen Spielbrett so bewegen, dass identische Objekte am Ende an denselben Koordinaten liegen. Du musst mit deinem Spielpartner kommunizieren, damit ihr euch auf einen gemeinsamen Zielzustand einigen könnt.} \\
\\ 
\\ 
\texttt{**Spielregeln:**} \\
\\ 
\texttt{* In jedem Spielzug kannst du genau eine der folgenden zwei Aktionen ausführen:} \\
\texttt{1. `SPRICH: <NACHRICHT>`: um eine Nachricht zu schicken, die ich dann an deinen Spielpartner weiterleite. Als Nachricht gilt alles bis zum nächsten Zeilenumbruch.} \\
\texttt{2. `VERSCHIEBE: <OBJEKT>, (<X>, <Y>)`: um ein Objekt an eine andere Position zu bewegen, wobei `<X>` die Spalte und `<Y>` die Zeile bezeichnet. Ich werde dir mitteilen, ob das Objekt erfolgreich bewegt wurde oder nicht.} \\
\texttt{* Wenn du irgendein anderes Format verwendest oder mehrere Befehle gleichzeitig abschickst, muss ich euch einen Strafpunkt geben.} \\
\texttt{* Wenn ihr gemeinsam mehr als 12 Strafpunkte ansammelt, verliert ihr das Spiel.} \\
\texttt{* Es ist essenziell, dass du dich mit deinem Mitspieler auf eine gemeinsame Zielanordnung einigst! Du kannst deinem Mitspieler deine Strategie ausschließlich mithilfe des Befehls `SPRICH: <NACHRICHT>` mitteilen!} \\
\\ 
\\ 
\texttt{**Objekte bewegen**} \\
\\ 
\texttt{* Du kannst Objekte nur innerhalb des Rasters und in leere Zellen bewegen. Leere Zellen erkennst am Zeichen '◌'.} \\
\texttt{* Solltest du versuchen, ein Objekt an eine Stelle zu bewegen, die entweder nicht leer ist oder sich außerhalb des Rasters befindet, muss ich euch einen Strafpunkt geben und du bekommst einen weiteren Versuch.} \\
\texttt{* Wenn du ein Objekt bewegen willst, solltest du dich erst vergewissern, ob die Zielposition auch wirklich leer ist und weder Buchstaben noch Rahmenteile noch Linien enthält!} \\
\\ 
\\ 
\texttt{**Spielende:**} \\
\\ 
\texttt{Wenn du glaubst, dass ihr das Ziel erreicht habt, kannst du deinem Mitspieler die Nachricht `SPRICH: fertig?` schicken. Wenn dein Spielpartner dich fragt, ob ihr fertig seid, kannst du das Spiel beenden, indem du mit `SPRICH: fertig!` antwortest.} \\
\\ 
\texttt{Ihr gewinnt gemeinsam, wenn ihr das Spiel innerhalb von \#max\_rounds Runden beendet. Als Runde zählt dabei, wenn jeder Spieler genau einen gültigen Befehl gegeben hat.} \\
\\ 
\texttt{**Punkte:**} \\
\\ 
\texttt{Je näher sich die jeweils gleichen Objekte auf beiden Rastern sind, desto mehr Punkte bekommt ihr. Für Strafpunkte gibt es Punktabzug. Könnt ihr den Rekord brechen?} \\
\\ 
\texttt{Schicke eine Nachricht an deinen Mitspieler um das Spiel zu starten!} \\
            }
        }
    }
    & \\ \\

    \theutterance \stepcounter{utterance}  
    & \multicolumn{4}{p{0.6\linewidth}}{
        \cellcolor[rgb]{0.9,0.9,0.9}{
            \makecell[{{p{\linewidth}}}]{
                \texttt{\tiny{[P1$\rangle$GM]}}
                \texttt{SPRICH: Hallo! Ich sehe meine Objekte an folgenden Positionen:} \\
\texttt{C: (1,1)} \\
\texttt{I: (7,1)} \\
\texttt{T: (6,1)} \\
\texttt{W: (4,2)} \\
\texttt{H: (3,6)} \\
\\ 
\texttt{Kannst du mir bitte deine Positionen nennen, damit wir einen gemeinsamen Zielzustand festlegen können?} \\
            }
        }
    }
    & & \\ \\

    \theutterance \stepcounter{utterance}  
    & & & \multicolumn{2}{p{0.3\linewidth}}{
        \cellcolor[rgb]{0.9,0.9,0.9}{
            \makecell[{{p{\linewidth}}}]{
                \texttt{\tiny{[GM$|$GM]}}
                \texttt{SPRICH: Hallo! Ich sehe meine Objekte an folgenden Positionen:} \\
\texttt{C: (1,1)} \\
\texttt{I: (7,1)} \\
\texttt{T: (6,1)} \\
\texttt{W: (4,2)} \\
\texttt{H: (3,6)} \\
\\ 
\texttt{Kannst du mir bitte deine Positionen nennen, damit wir einen gemeinsamen Zielzustand festlegen können?} \\
            }
        }
    }
    & & \\ \\

    \theutterance \stepcounter{utterance}  
    & & & \multicolumn{2}{p{0.3\linewidth}}{
        \cellcolor[rgb]{0.9,0.9,0.9}{
            \makecell[{{p{\linewidth}}}]{
                \texttt{\tiny{[GM$|$GM]}}
                \texttt{Invalid format: tail is not empty: '} \\
\texttt{C: (1,1)} \\
\texttt{I: (7,1)} \\
\texttt{T: (6,1)} \\
\texttt{W: (4,2)} \\
\texttt{H: (3,6)} \\
\\ 
\texttt{Kannst du mir bitte deine Positionen nennen, damit wir einen gemeinsamen Zielzustand festlegen können?} \\
            }
        }
    }
    & & \\ \\

    \theutterance \stepcounter{utterance}  
    & & \multicolumn{4}{p{0.6\linewidth}}{
        \cellcolor[rgb]{0.9,0.9,0.9}{
            \makecell[{{p{\linewidth}}}]{
                \texttt{\tiny{[P1$\langle$GM]}}
                \texttt{Strafpunkt: Deine Nachricht darf nichts nach dem Befehl enthalten! Stelle sicher, dass deine Antwort ausschließlich einen der beiden Befehle `SPRICH: <NACHRICHT>` und `VERSCHIEBE: <OBJEKT>, (<X>, <Y>)` und sonst nichts enthält!} \\
\\ 
\texttt{Ihr habt gemeinsam 1 von maximal 12 Strafpunkten angesammelt.} \\
\texttt{Bitte versuche es noch mal!} \\
            }
        }
    }
    & \\ \\

    \theutterance \stepcounter{utterance}  
    & \multicolumn{4}{p{0.6\linewidth}}{
        \cellcolor[rgb]{0.9,0.9,0.9}{
            \makecell[{{p{\linewidth}}}]{
                \texttt{\tiny{[P1$\rangle$GM]}}
                \texttt{SPRICH: Hallo! Ich sehe meine Objekte an folgenden Positionen: C: (1,1), I: (7,1), T: (6,1), W: (4,2), H: (3,6). Kannst du mir bitte deine Positionen nennen, damit wir einen gemeinsamen Zielzustand festlegen können?} \\
            }
        }
    }
    & & \\ \\

    \theutterance \stepcounter{utterance}  
    & & & \multicolumn{2}{p{0.3\linewidth}}{
        \cellcolor[rgb]{0.9,0.9,0.9}{
            \makecell[{{p{\linewidth}}}]{
                \texttt{\tiny{[GM$|$GM]}}
                \texttt{SPRICH: Hallo! Ich sehe meine Objekte an folgenden Positionen: C: (1,1), I: (7,1), T: (6,1), W: (4,2), H: (3,6). Kannst du mir bitte deine Positionen nennen, damit wir einen gemeinsamen Zielzustand festlegen können?} \\
            }
        }
    }
    & & \\ \\

    \theutterance \stepcounter{utterance}  
    & & \multicolumn{4}{p{0.6\linewidth}}{
        \cellcolor[rgb]{0.9,0.9,0.9}{
            \makecell[{{p{\linewidth}}}]{
                \texttt{\tiny{[GM$\rangle$P2]}}
                \texttt{Ich bin der Spielleiter, und du spielst ein kooperatives Spiel mit dem folgenden Raster als Spielbrett:} \\
\\ 
\texttt{```} \\
\texttt{ 1234567} \\
\texttt{╔═╤═╤═══╗} \\
\texttt{║H│◌│◌◌◌║ 1} \\
\texttt{║◌├─┤◌◌◌║ 2} \\
\texttt{╟─┤◌├─┬─╢ 3} \\
\texttt{║◌│◌│C│◌║ 4} \\
\texttt{╟─┼─┘◌│◌║ 5} \\
\texttt{║◌│◌◌┌┴─╢ 6} \\
\texttt{║T│WI│◌◌║ 7} \\
\texttt{╚═╧══╧══╝} \\
\\ 
\texttt{```} \\
\\ 
\texttt{* Am oberen Rand sind die x{-}Koordinaten angegeben, die nach rechts hin zunehmen, und am rechten Rand die y{-}Koordinaten, die nach unten zunehmen.} \\
\texttt{* Die folgenden Objekte sind zufällig auf deinem Raster platziert: 'W', 'I', 'T', 'C', 'H'.} \\
\\ 
\\ 
\texttt{Dein Mitspieler sieht eine Variation des Spielbretts, auf dem dieselben Objekte an anderen zufälligen Positionen platziert sind. Du kannst das Spielbrett des anderen Spielers nicht sehen, und er kann deins nicht sehen.} \\
\\ 
\\ 
\texttt{**Ziel des Spiels:**} \\
\\ 
\texttt{Beide Spieler müssen die Objekte auf ihrem jeweiligen Spielbrett so bewegen, dass identische Objekte am Ende an denselben Koordinaten liegen. Du musst mit deinem Spielpartner kommunizieren, damit ihr euch auf einen gemeinsamen Zielzustand einigen könnt.} \\
\\ 
\\ 
\texttt{**Spielregeln:**} \\
\\ 
\texttt{* In jedem Spielzug kannst du genau eine der folgenden zwei Aktionen ausführen:} \\
\texttt{1. `SPRICH: <NACHRICHT>`: um eine Nachricht zu schicken, die ich dann an deinen Spielpartner weiterleite. Als Nachricht gilt alles bis zum nächsten Zeilenumbruch.} \\
\texttt{2. `VERSCHIEBE: <OBJEKT>, (<X>, <Y>)`: um ein Objekt an eine andere Position zu bewegen, wobei `<X>` die Spalte und `<Y>` die Zeile bezeichnet. Ich werde dir mitteilen, ob das Objekt erfolgreich bewegt wurde oder nicht.} \\
\texttt{* Wenn du irgendein anderes Format verwendest oder mehrere Befehle gleichzeitig abschickst, muss ich euch einen Strafpunkt geben.} \\
\texttt{* Wenn ihr gemeinsam mehr als 12 Strafpunkte ansammelt, verliert ihr das Spiel.} \\
\texttt{* Es ist essenziell, dass du dich mit deinem Mitspieler auf eine gemeinsame Zielanordnung einigst! Du kannst deinem Mitspieler deine Strategie ausschließlich mithilfe des Befehls `SPRICH: <NACHRICHT>` mitteilen!} \\
\\ 
\\ 
\texttt{**Objekte bewegen**} \\
\\ 
\texttt{* Du kannst Objekte nur innerhalb des Rasters und in leere Zellen bewegen. Leere Zellen erkennst am Zeichen '◌'.} \\
\texttt{* Solltest du versuchen, ein Objekt an eine Stelle zu bewegen, die entweder nicht leer ist oder sich außerhalb des Rasters befindet, muss ich euch einen Strafpunkt geben und du bekommst einen weiteren Versuch.} \\
\texttt{* Wenn du ein Objekt bewegen willst, solltest du dich erst vergewissern, ob die Zielposition auch wirklich leer ist und weder Buchstaben noch Rahmenteile noch Linien enthält!} \\
\\ 
\\ 
\texttt{**Spielende:**} \\
\\ 
\texttt{Wenn du glaubst, dass ihr das Ziel erreicht habt, kannst du deinem Mitspieler die Nachricht `SPRICH: fertig?` schicken. Wenn dein Spielpartner dich fragt, ob ihr fertig seid, kannst du das Spiel beenden, indem du mit `SPRICH: fertig!` antwortest.} \\
\\ 
\texttt{Ihr gewinnt gemeinsam, wenn ihr das Spiel innerhalb von \#max\_rounds Runden beendet. Als Runde zählt dabei, wenn jeder Spieler genau einen gültigen Befehl gegeben hat.} \\
\\ 
\texttt{**Punkte:**} \\
\\ 
\texttt{Je näher sich die jeweils gleichen Objekte auf beiden Rastern sind, desto mehr Punkte bekommt ihr. Für Strafpunkte gibt es Punktabzug. Könnt ihr den Rekord brechen?} \\
\\ 
\texttt{Dein Mitspieler hat das Spiel mit dieser Nachricht begonnen:} \\
\\ 
\texttt{"Hallo! Ich sehe meine Objekte an folgenden Positionen: C: (1,1), I: (7,1), T: (6,1), W: (4,2), H: (3,6). Kannst du mir bitte deine Positionen nennen, damit wir einen gemeinsamen Zielzustand festlegen können?"} \\
\\ 
\texttt{Wie lautet dein erster Befehl?} \\
            }
        }
    }
    & \\ \\

    \theutterance \stepcounter{utterance}  
    & & & \multicolumn{4}{p{0.6\linewidth}}{
        \cellcolor[rgb]{0.9,0.9,0.9}{
            \makecell[{{p{\linewidth}}}]{
                \texttt{\tiny{[GM$\langle$P2]}}
                \texttt{SPRICH: Hallo! Meine Objekte befinden sich an folgenden Positionen: C: (3,4), H: (1,1), I: (4,7), T: (1,7), W: (4,7). Bitte beachte: Auf meinem Brett sind I und W beide in Zeile 7. I ist in Spalte 4, W ist in Spalte 4? Nein, warte. Lass mich das genau prüfen.} \\
\\ 
\texttt{Mein Raster:} \\
\texttt{Zeile 1: H an (1,1)} \\
\texttt{Zeile 4: C an (3,4)} \\
\texttt{Zeile 7: T an (1,7), W an (4,7), I an (4,7)?} \\
\\ 
\texttt{Schauen wir uns das Raster genau an:} \\
\texttt{```} \\
\texttt{ 1234567} \\
\texttt{╔═╤═╤═══╗} \\
\texttt{║H│◌│◌◌◌║ 1  {-}> H ist bei (1,1)} \\
\texttt{║◌├─┤◌◌◌║ 2} \\
\texttt{╟─┤◌├─┬─╢ 3} \\
\texttt{║◌│◌│C│◌║ 4  {-}> C ist bei (3,4)} \\
\texttt{╟─┼─┘◌│◌║ 5} \\
\texttt{║◌│◌◌┌┴─╢ 6} \\
\texttt{║T│WI│◌◌║ 7  {-}> T ist bei (1,7). W und I sind in den Zellen (4,7) und (5,7)?} \\
\texttt{╚═╧══╧══╝} \\
\texttt{```} \\
\\ 
\texttt{In Zeile 7 steht `T│WI│◌◌`.} \\
\texttt{Spalte 1: T} \\
\texttt{Spalte 2: | (Rahmen)} \\
\texttt{Spalte 3: W} \\
\texttt{Spalte 4: I} \\
\texttt{Spalte 5: | (Rahmen)} \\
\texttt{Spalte 6: ◌} \\
\texttt{Spalte 7: ◌} \\
\\ 
\texttt{Korrektur der Positionen auf meinem Brett:} \\
\texttt{H: (1,1)} \\
\texttt{C: (3,4)} \\
\texttt{T: (1,7)} \\
\texttt{W: (3,7)} \\
\texttt{I: (4,7)} \\
\\ 
\texttt{Deine Positionen:} \\
\texttt{C: (1,1)} \\
\texttt{I: (7,1)} \\
\texttt{T: (6,1)} \\
\texttt{W: (4,2)} \\
\texttt{H: (3,6)} \\
\\ 
\texttt{Lass uns als Zielzustand festlegen, dass wir alle Objekte auf die Positionen bringen, die ich gerade habe. Das ist für mich einfach zu koordinieren, da ich sie direkt sehe.} \\
\texttt{Zielzustand:} \\
\texttt{H: (1,1)} \\
\texttt{C: (3,4)} \\
\texttt{T: (1,7)} \\
\texttt{W: (3,7)} \\
\texttt{I: (4,7)} \\
\\ 
\texttt{Kannst du bestätigen, dass du diesen Zielzustand anstrebst? Wenn ja, fange ich an, meine Objekte zu bewegen, um sie in diese Positionen zu bringen, und du bewegst deine entsprechend.} \\
\\ 
\texttt{Zuerst muss ich H bewegen. H ist schon bei (1,1). Das ist gut.} \\
\texttt{C ist bei (3,4). Das ist gut.} \\
\texttt{T ist bei (1,7). Das ist gut.} \\
\texttt{W ist bei (3,7). Das ist gut.} \\
\texttt{I ist bei (4,7). Das ist gut.} \\
\\ 
\texttt{Warte, ich habe meine Positionen bereits korrekt identifiziert. Alle meine Objekte sind bereits an den Positionen, die ich als Ziel definiert habe. Das bedeutet, ich muss nichts verschieben, um meinen Teil des Zielzustands zu erreichen.} \\
\\ 
\texttt{Allerdings muss ich sicherstellen, dass die Zielpositionen auf *deinem* Brett auch frei sind, damit du deine Objekte dorthin bewegen kannst. Da ich dein Brett nicht sehe, ist das riskant.} \\
\\ 
\texttt{Eine bessere Strategie wäre es, einen neutralen Punkt zu wählen oder einfach zu versuchen, die Objekte auf eine Standardposition zu bringen, z.B. die linke obere Ecke, aber das Raster ist klein.} \\
\\ 
\texttt{Lass uns einen anderen Ansatz wählen. Wir einigen uns darauf, dass wir alle Objekte auf die Positionen bewegen, die ich *jetzt* habe. Da ich meine Objekte bereits an diesen Positionen habe, muss ich nur noch prüfen, ob diese Positionen auf deinem Brett frei sind. Das kann ich nicht wissen.} \\
\\ 
\texttt{Stattdessen: Ich bewege meine Objekte nicht, sondern warte ab, was du tust? Nein, ich muss aktiv werden.} \\
\\ 
\texttt{Lass uns einen einfachen Zielzustand wählen, der wahrscheinlich auf beiden Brettern erreichbar ist. Zum Beispiel:} \\
\texttt{H {-}> (1,1)} \\
\texttt{C {-}> (1,2)} \\
\texttt{T {-}> (1,3)} \\
\texttt{W {-}> (1,4)} \\
\texttt{I {-}> (1,5)} \\
\\ 
\texttt{Das ist eine Spalte. Aber ich weiß nicht, ob diese Zellen bei dir frei sind.} \\
\\ 
\texttt{Am sichersten ist es, wenn wir uns auf die Positionen einigen, die ich *aktuell* habe, und du prüfst, ob diese bei dir frei sind. Wenn nicht, müssen wir einen anderen Plan schmieden.} \\
\\ 
\texttt{Ich sage dir also meine aktuellen Positionen noch einmal klar und deutlich und schlage vor, dass wir als Ziel genau diese nehmen.} \\
\\ 
\texttt{Meine Positionen:} \\
\texttt{H: (1,1)} \\
\texttt{C: (3,4)} \\
\texttt{T: (1,7)} \\
\texttt{W: (3,7)} \\
\texttt{I: (4,7)} \\
\\ 
\texttt{Bitte bestätige, ob du diese Positionen als Ziel annehmen kannst. Wenn ja, bewege ich keine Objekte, da sie bereits dort sind, und du bewegst deine.} \\
\\ 
\texttt{Da ich aber einen Befehl abgeben muss und ich keine Objekte verschieben muss (sie sind schon am Ziel), kann ich theoretisch `VERSCHIEBE` nicht nutzen, es sei denn, ich verschiebe ein Objekt kurz und zurück, was Zeit kostet. Oder ich sende nur eine Nachricht.} \\
\\ 
\texttt{Die Regel sagt: "In jedem Spielzug kannst du genau eine der folgenden zwei Aktionen ausführen". Eine Nachricht ist eine Aktion.} \\
\\ 
\texttt{Also sende ich die Nachricht mit den Positionen und dem Vorschlag.} \\
\\ 
\texttt{SPRICH: Hallo! Meine Objekte sind: H:(1,1), C:(3,4), T:(1,7), W:(3,7), I:(4,7). Ich schlage vor, wir nehmen diese Positionen als gemeinsamen Zielzustand. Bitte bestätige, ob diese Positionen auf deinem Brett frei sind, damit du deine Objekte dorthin bewegen kannst. Wenn ja, bewege ich nichts, da meine Objekte bereits dort sind.} \\
            }
        }
    }
     \\ \\

    \theutterance \stepcounter{utterance}  
    & & & \multicolumn{2}{p{0.3\linewidth}}{
        \cellcolor[rgb]{0.9,0.9,0.9}{
            \makecell[{{p{\linewidth}}}]{
                \texttt{\tiny{[GM$|$GM]}}
                \texttt{SPRICH: Hallo! Meine Objekte befinden sich an folgenden Positionen: C: (3,4), H: (1,1), I: (4,7), T: (1,7), W: (4,7). Bitte beachte: Auf meinem Brett sind I und W beide in Zeile 7. I ist in Spalte 4, W ist in Spalte 4? Nein, warte. Lass mich das genau prüfen.} \\
\\ 
\texttt{Mein Raster:} \\
\texttt{Zeile 1: H an (1,1)} \\
\texttt{Zeile 4: C an (3,4)} \\
\texttt{Zeile 7: T an (1,7), W an (4,7), I an (4,7)?} \\
\\ 
\texttt{Schauen wir uns das Raster genau an:} \\
\texttt{```} \\
\texttt{ 1234567} \\
\texttt{╔═╤═╤═══╗} \\
\texttt{║H│◌│◌◌◌║ 1  {-}> H ist bei (1,1)} \\
\texttt{║◌├─┤◌◌◌║ 2} \\
\texttt{╟─┤◌├─┬─╢ 3} \\
\texttt{║◌│◌│C│◌║ 4  {-}> C ist bei (3,4)} \\
\texttt{╟─┼─┘◌│◌║ 5} \\
\texttt{║◌│◌◌┌┴─╢ 6} \\
\texttt{║T│WI│◌◌║ 7  {-}> T ist bei (1,7). W und I sind in den Zellen (4,7) und (5,7)?} \\
\texttt{╚═╧══╧══╝} \\
\texttt{```} \\
\\ 
\texttt{In Zeile 7 steht `T│WI│◌◌`.} \\
\texttt{Spalte 1: T} \\
\texttt{Spalte 2: | (Rahmen)} \\
\texttt{Spalte 3: W} \\
\texttt{Spalte 4: I} \\
\texttt{Spalte 5: | (Rahmen)} \\
\texttt{Spalte 6: ◌} \\
\texttt{Spalte 7: ◌} \\
\\ 
\texttt{Korrektur der Positionen auf meinem Brett:} \\
\texttt{H: (1,1)} \\
\texttt{C: (3,4)} \\
\texttt{T: (1,7)} \\
\texttt{W: (3,7)} \\
\texttt{I: (4,7)} \\
\\ 
\texttt{Deine Positionen:} \\
\texttt{C: (1,1)} \\
\texttt{I: (7,1)} \\
\texttt{T: (6,1)} \\
\texttt{W: (4,2)} \\
\texttt{H: (3,6)} \\
\\ 
\texttt{Lass uns als Zielzustand festlegen, dass wir alle Objekte auf die Positionen bringen, die ich gerade habe. Das ist für mich einfach zu koordinieren, da ich sie direkt sehe.} \\
\texttt{Zielzustand:} \\
\texttt{H: (1,1)} \\
\texttt{C: (3,4)} \\
\texttt{T: (1,7)} \\
\texttt{W: (3,7)} \\
\texttt{I: (4,7)} \\
\\ 
\texttt{Kannst du bestätigen, dass du diesen Zielzustand anstrebst? Wenn ja, fange ich an, meine Objekte zu bewegen, um sie in diese Positionen zu bringen, und du bewegst deine entsprechend.} \\
\\ 
\texttt{Zuerst muss ich H bewegen. H ist schon bei (1,1). Das ist gut.} \\
\texttt{C ist bei (3,4). Das ist gut.} \\
\texttt{T ist bei (1,7). Das ist gut.} \\
\texttt{W ist bei (3,7). Das ist gut.} \\
\texttt{I ist bei (4,7). Das ist gut.} \\
\\ 
\texttt{Warte, ich habe meine Positionen bereits korrekt identifiziert. Alle meine Objekte sind bereits an den Positionen, die ich als Ziel definiert habe. Das bedeutet, ich muss nichts verschieben, um meinen Teil des Zielzustands zu erreichen.} \\
\\ 
\texttt{Allerdings muss ich sicherstellen, dass die Zielpositionen auf *deinem* Brett auch frei sind, damit du deine Objekte dorthin bewegen kannst. Da ich dein Brett nicht sehe, ist das riskant.} \\
\\ 
\texttt{Eine bessere Strategie wäre es, einen neutralen Punkt zu wählen oder einfach zu versuchen, die Objekte auf eine Standardposition zu bringen, z.B. die linke obere Ecke, aber das Raster ist klein.} \\
\\ 
\texttt{Lass uns einen anderen Ansatz wählen. Wir einigen uns darauf, dass wir alle Objekte auf die Positionen bewegen, die ich *jetzt* habe. Da ich meine Objekte bereits an diesen Positionen habe, muss ich nur noch prüfen, ob diese Positionen auf deinem Brett frei sind. Das kann ich nicht wissen.} \\
\\ 
\texttt{Stattdessen: Ich bewege meine Objekte nicht, sondern warte ab, was du tust? Nein, ich muss aktiv werden.} \\
\\ 
\texttt{Lass uns einen einfachen Zielzustand wählen, der wahrscheinlich auf beiden Brettern erreichbar ist. Zum Beispiel:} \\
\texttt{H {-}> (1,1)} \\
\texttt{C {-}> (1,2)} \\
\texttt{T {-}> (1,3)} \\
\texttt{W {-}> (1,4)} \\
\texttt{I {-}> (1,5)} \\
\\ 
\texttt{Das ist eine Spalte. Aber ich weiß nicht, ob diese Zellen bei dir frei sind.} \\
\\ 
\texttt{Am sichersten ist es, wenn wir uns auf die Positionen einigen, die ich *aktuell* habe, und du prüfst, ob diese bei dir frei sind. Wenn nicht, müssen wir einen anderen Plan schmieden.} \\
\\ 
\texttt{Ich sage dir also meine aktuellen Positionen noch einmal klar und deutlich und schlage vor, dass wir als Ziel genau diese nehmen.} \\
\\ 
\texttt{Meine Positionen:} \\
\texttt{H: (1,1)} \\
\texttt{C: (3,4)} \\
\texttt{T: (1,7)} \\
\texttt{W: (3,7)} \\
\texttt{I: (4,7)} \\
\\ 
\texttt{Bitte bestätige, ob du diese Positionen als Ziel annehmen kannst. Wenn ja, bewege ich keine Objekte, da sie bereits dort sind, und du bewegst deine.} \\
\\ 
\texttt{Da ich aber einen Befehl abgeben muss und ich keine Objekte verschieben muss (sie sind schon am Ziel), kann ich theoretisch `VERSCHIEBE` nicht nutzen, es sei denn, ich verschiebe ein Objekt kurz und zurück, was Zeit kostet. Oder ich sende nur eine Nachricht.} \\
\\ 
\texttt{Die Regel sagt: "In jedem Spielzug kannst du genau eine der folgenden zwei Aktionen ausführen". Eine Nachricht ist eine Aktion.} \\
\\ 
\texttt{Also sende ich die Nachricht mit den Positionen und dem Vorschlag.} \\
\\ 
\texttt{SPRICH: Hallo! Meine Objekte sind: H:(1,1), C:(3,4), T:(1,7), W:(3,7), I:(4,7). Ich schlage vor, wir nehmen diese Positionen als gemeinsamen Zielzustand. Bitte bestätige, ob diese Positionen auf deinem Brett frei sind, damit du deine Objekte dorthin bewegen kannst. Wenn ja, bewege ich nichts, da meine Objekte bereits dort sind.} \\
            }
        }
    }
    & & \\ \\

    \theutterance \stepcounter{utterance}  
    & & & \multicolumn{2}{p{0.3\linewidth}}{
        \cellcolor[rgb]{0.9,0.9,0.9}{
            \makecell[{{p{\linewidth}}}]{
                \texttt{\tiny{[GM$|$GM]}}
                \texttt{Invalid format: head is not empty: 'SPRICH: Hallo! Meine Objekte befinden sich an folgenden Positionen: C: (3,4), H: (1,1), I: (4,7), T: (1,7), W: (4,7). Bitte beachte: Auf meinem Brett sind I und W beide in Zeile 7. I ist in Spalte 4, W ist in Spalte 4? Nein, warte. Lass mich das genau prüfen.} \\
\\ 
\texttt{Mein Raster:} \\
\texttt{Zeile 1: H an (1,1)} \\
\texttt{Zeile 4: C an (3,4)} \\
\texttt{Zeile 7: T an (1,7), W an (4,7), I an (4,7)?} \\
\\ 
\texttt{Schauen wir uns das Raster genau an:} \\
\\ 
\texttt{ 1234567} \\
\texttt{╔═╤═╤═══╗} \\
\texttt{║H│◌│◌◌◌║ 1  {-}> H ist bei (1,1)} \\
\texttt{║◌├─┤◌◌◌║ 2} \\
\texttt{╟─┤◌├─┬─╢ 3} \\
\texttt{║◌│◌│C│◌║ 4  {-}> C ist bei (3,4)} \\
\texttt{╟─┼─┘◌│◌║ 5} \\
\texttt{║◌│◌◌┌┴─╢ 6} \\
\texttt{║T│WI│◌◌║ 7  {-}> T ist bei (1,7). W und I sind in den Zellen (4,7) und (5,7)?} \\
\texttt{╚═╧══╧══╝} \\
\\ 
\\ 
\texttt{In Zeile 7 steht T│WI│◌◌.} \\
\texttt{Spalte 1: T} \\
\texttt{Spalte 2: | (Rahmen)} \\
\texttt{Spalte 3: W} \\
\texttt{Spalte 4: I} \\
\texttt{Spalte 5: | (Rahmen)} \\
\texttt{Spalte 6: ◌} \\
\texttt{Spalte 7: ◌} \\
\\ 
\texttt{Korrektur der Positionen auf meinem Brett:} \\
\texttt{H: (1,1)} \\
\texttt{C: (3,4)} \\
\texttt{T: (1,7)} \\
\texttt{W: (3,7)} \\
\texttt{I: (4,7)} \\
\\ 
\texttt{Deine Positionen:} \\
\texttt{C: (1,1)} \\
\texttt{I: (7,1)} \\
\texttt{T: (6,1)} \\
\texttt{W: (4,2)} \\
\texttt{H: (3,6)} \\
\\ 
\texttt{Lass uns als Zielzustand festlegen, dass wir alle Objekte auf die Positionen bringen, die ich gerade habe. Das ist für mich einfach zu koordinieren, da ich sie direkt sehe.} \\
\texttt{Zielzustand:} \\
\texttt{H: (1,1)} \\
\texttt{C: (3,4)} \\
\texttt{T: (1,7)} \\
\texttt{W: (3,7)} \\
\texttt{I: (4,7)} \\
\\ 
\texttt{Kannst du bestätigen, dass du diesen Zielzustand anstrebst? Wenn ja, fange ich an, meine Objekte zu bewegen, um sie in diese Positionen zu bringen, und du bewegst deine entsprechend.} \\
\\ 
\texttt{Zuerst muss ich H bewegen. H ist schon bei (1,1). Das ist gut.} \\
\texttt{C ist bei (3,4). Das ist gut.} \\
\texttt{T ist bei (1,7). Das ist gut.} \\
\texttt{W ist bei (3,7). Das ist gut.} \\
\texttt{I ist bei (4,7). Das ist gut.} \\
\\ 
\texttt{Warte, ich habe meine Positionen bereits korrekt identifiziert. Alle meine Objekte sind bereits an den Positionen, die ich als Ziel definiert habe. Das bedeutet, ich muss nichts verschieben, um meinen Teil des Zielzustands zu erreichen.} \\
\\ 
\texttt{Allerdings muss ich sicherstellen, dass die Zielpositionen auf *deinem* Brett auch frei sind, damit du deine Objekte dorthin bewegen kannst. Da ich dein Brett nicht sehe, ist das riskant.} \\
\\ 
\texttt{Eine bessere Strategie wäre es, einen neutralen Punkt zu wählen oder einfach zu versuchen, die Objekte auf eine Standardposition zu bringen, z.B. die linke obere Ecke, aber das Raster ist klein.} \\
\\ 
\texttt{Lass uns einen anderen Ansatz wählen. Wir einigen uns darauf, dass wir alle Objekte auf die Positionen bewegen, die ich *jetzt* habe. Da ich meine Objekte bereits an diesen Positionen habe, muss ich nur noch prüfen, ob diese Positionen auf deinem Brett frei sind. Das kann ich nicht wissen.} \\
\\ 
\texttt{Stattdessen: Ich bewege meine Objekte nicht, sondern warte ab, was du tust? Nein, ich muss aktiv werden.} \\
\\ 
\texttt{Lass uns einen einfachen Zielzustand wählen, der wahrscheinlich auf beiden Brettern erreichbar ist. Zum Beispiel:} \\
\texttt{H {-}> (1,1)} \\
\texttt{C {-}> (1,2)} \\
\texttt{T {-}> (1,3)} \\
\texttt{W {-}> (1,4)} \\
\texttt{I {-}> (1,5)} \\
\\ 
\texttt{Das ist eine Spalte. Aber ich weiß nicht, ob diese Zellen bei dir frei sind.} \\
\\ 
\texttt{Am sichersten ist es, wenn wir uns auf die Positionen einigen, die ich *aktuell* habe, und du prüfst, ob diese bei dir frei sind. Wenn nicht, müssen wir einen anderen Plan schmieden.} \\
\\ 
\texttt{Ich sage dir also meine aktuellen Positionen noch einmal klar und deutlich und schlage vor, dass wir als Ziel genau diese nehmen.} \\
\\ 
\texttt{Meine Positionen:} \\
\texttt{H: (1,1)} \\
\texttt{C: (3,4)} \\
\texttt{T: (1,7)} \\
\texttt{W: (3,7)} \\
\texttt{I: (4,7)} \\
\\ 
\texttt{Bitte bestätige, ob du diese Positionen als Ziel annehmen kannst. Wenn ja, bewege ich keine Objekte, da sie bereits dort sind, und du bewegst deine.} \\
\\ 
\texttt{Da ich aber einen Befehl abgeben muss und ich keine Objekte verschieben muss (sie sind schon am Ziel), kann ich theoretisch VERSCHIEBE nicht nutzen, es sei denn, ich verschiebe ein Objekt kurz und zurück, was Zeit kostet. Oder ich sende nur eine Nachricht.} \\
\\ 
\texttt{Die Regel sagt: "In jedem Spielzug kannst du genau eine der folgenden zwei Aktionen ausführen". Eine Nachricht ist eine Aktion.} \\
\\ 
\texttt{Also sende ich die Nachricht mit den Positionen und dem Vorschlag.} \\
\\ 
\texttt{'} \\
            }
        }
    }
    & & \\ \\

    \theutterance \stepcounter{utterance}  
    & & \multicolumn{4}{p{0.6\linewidth}}{
        \cellcolor[rgb]{0.9,0.9,0.9}{
            \makecell[{{p{\linewidth}}}]{
                \texttt{\tiny{[GM$\rangle$P2]}}
                \texttt{Strafpunkt: Deine Nachricht darf nichts vor dem Befehl enthalten! Stelle sicher, dass deine Antwort ausschließlich einen der beiden Befehle `SPRICH: <NACHRICHT>` und `VERSCHIEBE: <OBJEKT>, (<X>, <Y>)` und sonst nichts enthält!} \\
\\ 
\texttt{Ihr habt gemeinsam 2 von maximal 12 Strafpunkten angesammelt.} \\
\texttt{Bitte versuche es noch mal!} \\
            }
        }
    }
    & \\ \\

    \theutterance \stepcounter{utterance}  
    & & & \multicolumn{4}{p{0.6\linewidth}}{
        \cellcolor[rgb]{0.9,0.9,0.9}{
            \makecell[{{p{\linewidth}}}]{
                \texttt{\tiny{[GM$\langle$P2]}}
                \texttt{SPRICH: Hallo! Meine Objekte sind: H:(1,1), C:(3,4), T:(1,7), W:(3,7), I:(4,7). Ich schlage vor, wir nehmen diese Positionen als gemeinsamen Zielzustand. Bitte bestätige, ob diese Positionen auf deinem Brett frei sind, damit du deine Objekte dorthin bewegen kannst.} \\
            }
        }
    }
     \\ \\

    \theutterance \stepcounter{utterance}  
    & & & \multicolumn{2}{p{0.3\linewidth}}{
        \cellcolor[rgb]{0.9,0.9,0.9}{
            \makecell[{{p{\linewidth}}}]{
                \texttt{\tiny{[GM$|$GM]}}
                \texttt{SPRICH: Hallo! Meine Objekte sind: H:(1,1), C:(3,4), T:(1,7), W:(3,7), I:(4,7). Ich schlage vor, wir nehmen diese Positionen als gemeinsamen Zielzustand. Bitte bestätige, ob diese Positionen auf deinem Brett frei sind, damit du deine Objekte dorthin bewegen kannst.} \\
            }
        }
    }
    & & \\ \\

    \theutterance \stepcounter{utterance}  
    & & \multicolumn{4}{p{0.6\linewidth}}{
        \cellcolor[rgb]{0.9,0.9,0.9}{
            \makecell[{{p{\linewidth}}}]{
                \texttt{\tiny{[P1$\langle$GM]}}
                \texttt{Deine Nachricht wurde an deinen Mitspieler übermittelt.} \\
\\ 
\\ 
\texttt{Ihr seid gerade bei Runde 1 von maximal 20.} \\
\texttt{Ihr habt gemeinsam 2 von maximal 12 Strafpunkten angesammelt.} \\
\texttt{Dein Mitspieler hat dir diese Nachricht geschickt:} \\
\\ 
\\ 
\texttt{"Hallo! Meine Objekte sind: H:(1,1), C:(3,4), T:(1,7), W:(3,7), I:(4,7). Ich schlage vor, wir nehmen diese Positionen als gemeinsamen Zielzustand. Bitte bestätige, ob diese Positionen auf deinem Brett frei sind, damit du deine Objekte dorthin bewegen kannst."} \\
\\ 
\\ 
\texttt{Wie lautet dein nächster Befehl?} \\
            }
        }
    }
    & \\ \\

    \theutterance \stepcounter{utterance}  
    & \multicolumn{4}{p{0.6\linewidth}}{
        \cellcolor[rgb]{0.9,0.9,0.9}{
            \makecell[{{p{\linewidth}}}]{
                \texttt{\tiny{[P1$\rangle$GM]}}
                \texttt{SPRICH: Ich überprüfe die Zielpositionen auf meinem Brett:} \\
\texttt{H:(1,1) {-}> Besetzt mit C} \\
\texttt{C:(3,4) {-}> Frei} \\
\texttt{T:(1,7) {-}> Frei} \\
\texttt{W:(3,7) {-}> Frei} \\
\texttt{I:(4,7) {-}> Frei} \\
\\ 
\texttt{Da H:(1,1) bei mir bereits von C besetzt ist, können wir diesen Zielzustand nicht direkt so übernehmen. Ich schlage vor, wir tauschen die Positionen von C und H.} \\
\texttt{Neuer Vorschlag für den Zielzustand:} \\
\texttt{C:(3,6) {[}wo H ist{]}} \\
\texttt{H:(1,1) {[}wo C ist{]}} \\
\texttt{T:(1,7)} \\
\texttt{W:(3,7)} \\
\texttt{I:(4,7)} \\
\\ 
\texttt{Bitte prüfe, ob diese neuen Positionen auf deinem Brett frei sind.} \\
            }
        }
    }
    & & \\ \\

    \theutterance \stepcounter{utterance}  
    & & & \multicolumn{2}{p{0.3\linewidth}}{
        \cellcolor[rgb]{0.9,0.9,0.9}{
            \makecell[{{p{\linewidth}}}]{
                \texttt{\tiny{[GM$|$GM]}}
                \texttt{SPRICH: Ich überprüfe die Zielpositionen auf meinem Brett:} \\
\texttt{H:(1,1) {-}> Besetzt mit C} \\
\texttt{C:(3,4) {-}> Frei} \\
\texttt{T:(1,7) {-}> Frei} \\
\texttt{W:(3,7) {-}> Frei} \\
\texttt{I:(4,7) {-}> Frei} \\
\\ 
\texttt{Da H:(1,1) bei mir bereits von C besetzt ist, können wir diesen Zielzustand nicht direkt so übernehmen. Ich schlage vor, wir tauschen die Positionen von C und H.} \\
\texttt{Neuer Vorschlag für den Zielzustand:} \\
\texttt{C:(3,6) {[}wo H ist{]}} \\
\texttt{H:(1,1) {[}wo C ist{]}} \\
\texttt{T:(1,7)} \\
\texttt{W:(3,7)} \\
\texttt{I:(4,7)} \\
\\ 
\texttt{Bitte prüfe, ob diese neuen Positionen auf deinem Brett frei sind.} \\
            }
        }
    }
    & & \\ \\

    \theutterance \stepcounter{utterance}  
    & & & \multicolumn{2}{p{0.3\linewidth}}{
        \cellcolor[rgb]{0.9,0.9,0.9}{
            \makecell[{{p{\linewidth}}}]{
                \texttt{\tiny{[GM$|$GM]}}
                \texttt{Invalid format: tail is not empty: '} \\
\texttt{H:(1,1) {-}> Besetzt mit C} \\
\texttt{C:(3,4) {-}> Frei} \\
\texttt{T:(1,7) {-}> Frei} \\
\texttt{W:(3,7) {-}> Frei} \\
\texttt{I:(4,7) {-}> Frei} \\
\\ 
\texttt{Da H:(1,1) bei mir bereits von C besetzt ist, können wir diesen Zielzustand nicht direkt so übernehmen. Ich schlage vor, wir tauschen die Positionen von C und H.} \\
\texttt{Neuer Vorschlag für den Zielzustand:} \\
\texttt{C:(3,6) {[}wo H ist{]}} \\
\texttt{H:(1,1) {[}wo C ist{]}} \\
\texttt{T:(1,7)} \\
\texttt{W:(3,7)} \\
\texttt{I:(4,7)} \\
\\ 
\texttt{Bitte prüfe, ob diese neuen Positionen auf deinem Brett frei sind.} \\
            }
        }
    }
    & & \\ \\

    \theutterance \stepcounter{utterance}  
    & & \multicolumn{4}{p{0.6\linewidth}}{
        \cellcolor[rgb]{0.9,0.9,0.9}{
            \makecell[{{p{\linewidth}}}]{
                \texttt{\tiny{[P1$\langle$GM]}}
                \texttt{Strafpunkt: Deine Nachricht darf nichts nach dem Befehl enthalten! Stelle sicher, dass deine Antwort ausschließlich einen der beiden Befehle `SPRICH: <NACHRICHT>` und `VERSCHIEBE: <OBJEKT>, (<X>, <Y>)` und sonst nichts enthält!} \\
\\ 
\texttt{Ihr habt gemeinsam 3 von maximal 12 Strafpunkten angesammelt.} \\
\texttt{Bitte versuche es noch mal!} \\
            }
        }
    }
    & \\ \\

    \theutterance \stepcounter{utterance}  
    & \multicolumn{4}{p{0.6\linewidth}}{
        \cellcolor[rgb]{0.9,0.9,0.9}{
            \makecell[{{p{\linewidth}}}]{
                \texttt{\tiny{[P1$\rangle$GM]}}
                \texttt{SPRICH: Ich überprüfe die Zielpositionen auf meinem Brett: H:(1,1) ist besetzt mit C. C:(3,4) ist frei. T:(1,7) ist frei. W:(3,7) ist frei. I:(4,7) ist frei. Da H:(1,1) besetzt ist, schlage ich vor, wir tauschen C und H. Neuer Zielzustand: C:(3,6), H:(1,1), T:(1,7), W:(3,7), I:(4,7). Bitte prüfe, ob diese neuen Positionen auf deinem Brett frei sind.} \\
            }
        }
    }
    & & \\ \\

    \theutterance \stepcounter{utterance}  
    & & & \multicolumn{2}{p{0.3\linewidth}}{
        \cellcolor[rgb]{0.9,0.9,0.9}{
            \makecell[{{p{\linewidth}}}]{
                \texttt{\tiny{[GM$|$GM]}}
                \texttt{SPRICH: Ich überprüfe die Zielpositionen auf meinem Brett: H:(1,1) ist besetzt mit C. C:(3,4) ist frei. T:(1,7) ist frei. W:(3,7) ist frei. I:(4,7) ist frei. Da H:(1,1) besetzt ist, schlage ich vor, wir tauschen C und H. Neuer Zielzustand: C:(3,6), H:(1,1), T:(1,7), W:(3,7), I:(4,7). Bitte prüfe, ob diese neuen Positionen auf deinem Brett frei sind.} \\
            }
        }
    }
    & & \\ \\

    \theutterance \stepcounter{utterance}  
    & & \multicolumn{4}{p{0.6\linewidth}}{
        \cellcolor[rgb]{0.9,0.9,0.9}{
            \makecell[{{p{\linewidth}}}]{
                \texttt{\tiny{[GM$\rangle$P2]}}
                \texttt{Deine Nachricht wurde an deinen Mitspieler übermittelt.} \\
\\ 
\\ 
\texttt{Ihr seid gerade bei Runde 2 von maximal 20.} \\
\texttt{Ihr habt gemeinsam 3 von maximal 12 Strafpunkten angesammelt.} \\
\texttt{Dein Mitspieler hat dir diese Nachricht geschickt:} \\
\\ 
\\ 
\texttt{"Ich überprüfe die Zielpositionen auf meinem Brett: H:(1,1) ist besetzt mit C. C:(3,4) ist frei. T:(1,7) ist frei. W:(3,7) ist frei. I:(4,7) ist frei. Da H:(1,1) besetzt ist, schlage ich vor, wir tauschen C und H. Neuer Zielzustand: C:(3,6), H:(1,1), T:(1,7), W:(3,7), I:(4,7). Bitte prüfe, ob diese neuen Positionen auf deinem Brett frei sind."} \\
\\ 
\\ 
\texttt{Wie lautet dein nächster Befehl?} \\
            }
        }
    }
    & \\ \\

    \theutterance \stepcounter{utterance}  
    & & & \multicolumn{4}{p{0.6\linewidth}}{
        \cellcolor[rgb]{0.9,0.9,0.9}{
            \makecell[{{p{\linewidth}}}]{
                \texttt{\tiny{[GM$\langle$P2]}}
                \texttt{SPRICH: Ich prüfe die neuen Zielpositionen auf meinem Brett: C:(3,6) ist leer. H:(1,1) ist besetzt mit H. T:(1,7) ist besetzt mit T. W:(3,7) ist besetzt mit W. I:(4,7) ist besetzt mit I. Da H, T, W und I bereits an ihren Zielpositionen sind, aber C nicht, und H an (1,1) ist, müssen wir einen anderen Plan finden. Da H, T, W, I bereits an den Zielen sind, die wir vereinbart hatten (außer dass du C verschieben willst), und ich C an (3,4) habe, ist (3,6) bei mir leer. Aber H ist bei mir an (1,1), was dein Ziel für H ist. Also ist (1,1) bei mir besetzt. Das bedeutet, wir können H nicht nach (1,1) bewegen, wenn es dort schon ist? Nein, H ist schon da. Das Problem ist, dass du C nach (3,6) bewegen willst, aber bei mir ist (3,6) leer. Das ist gut für mich. Aber du willst H nach (1,1). Bei mir ist H schon an (1,1). Das ist auch gut. Aber du hast gesagt, H:(1,1) ist bei dir besetzt mit C. Also musst du C weg bewegen. Du willst C nach (3,6). Bei mir ist (3,6) leer. Das passt. Aber was ist mit den anderen? T, W, I sind bei dir an (1,7), (3,7), (4,7). Bei mir sind sie auch an (1,7), (3,7), (4,7). Also sind T, W, I bereits an den Zielen. C soll von (1,1) zu (3,6). Bei mir ist C an (3,4). Ich muss C von (3,4) nach (3,6) bewegen? Nein, mein Ziel ist es, dass meine Objekte an denselben Koordinaten wie deine sind. Wenn dein Ziel C:(3,6) ist, dann muss ich C nach (3,6) bewegen. Ist (3,6) bei mir leer? Ja. Also bewege ich C nach (3,6). Aber warte, du hast gesagt "Neuer Zielzustand: C:(3,6), H:(1,1), T:(1,7), W:(3,7), I:(4,7)". Das bedeutet, du willst, dass dein C an (3,6) ist und mein C auch an (3,6). Und dein H an (1,1) und mein H an (1,1). Da mein H schon an (1,1) ist, ist das gut. Da mein T, W, I schon an (1,7), (3,7), (4,7) sind, ist das gut. Ich muss nur mein C von (3,4) nach (3,6) bewegen. Ist (3,6) bei mir leer? Ja. Also verschiebe ich C nach (3,6).} \\
            }
        }
    }
     \\ \\

    \theutterance \stepcounter{utterance}  
    & & & \multicolumn{2}{p{0.3\linewidth}}{
        \cellcolor[rgb]{0.9,0.9,0.9}{
            \makecell[{{p{\linewidth}}}]{
                \texttt{\tiny{[GM$|$GM]}}
                \texttt{SPRICH: Ich prüfe die neuen Zielpositionen auf meinem Brett: C:(3,6) ist leer. H:(1,1) ist besetzt mit H. T:(1,7) ist besetzt mit T. W:(3,7) ist besetzt mit W. I:(4,7) ist besetzt mit I. Da H, T, W und I bereits an ihren Zielpositionen sind, aber C nicht, und H an (1,1) ist, müssen wir einen anderen Plan finden. Da H, T, W, I bereits an den Zielen sind, die wir vereinbart hatten (außer dass du C verschieben willst), und ich C an (3,4) habe, ist (3,6) bei mir leer. Aber H ist bei mir an (1,1), was dein Ziel für H ist. Also ist (1,1) bei mir besetzt. Das bedeutet, wir können H nicht nach (1,1) bewegen, wenn es dort schon ist? Nein, H ist schon da. Das Problem ist, dass du C nach (3,6) bewegen willst, aber bei mir ist (3,6) leer. Das ist gut für mich. Aber du willst H nach (1,1). Bei mir ist H schon an (1,1). Das ist auch gut. Aber du hast gesagt, H:(1,1) ist bei dir besetzt mit C. Also musst du C weg bewegen. Du willst C nach (3,6). Bei mir ist (3,6) leer. Das passt. Aber was ist mit den anderen? T, W, I sind bei dir an (1,7), (3,7), (4,7). Bei mir sind sie auch an (1,7), (3,7), (4,7). Also sind T, W, I bereits an den Zielen. C soll von (1,1) zu (3,6). Bei mir ist C an (3,4). Ich muss C von (3,4) nach (3,6) bewegen? Nein, mein Ziel ist es, dass meine Objekte an denselben Koordinaten wie deine sind. Wenn dein Ziel C:(3,6) ist, dann muss ich C nach (3,6) bewegen. Ist (3,6) bei mir leer? Ja. Also bewege ich C nach (3,6). Aber warte, du hast gesagt "Neuer Zielzustand: C:(3,6), H:(1,1), T:(1,7), W:(3,7), I:(4,7)". Das bedeutet, du willst, dass dein C an (3,6) ist und mein C auch an (3,6). Und dein H an (1,1) und mein H an (1,1). Da mein H schon an (1,1) ist, ist das gut. Da mein T, W, I schon an (1,7), (3,7), (4,7) sind, ist das gut. Ich muss nur mein C von (3,4) nach (3,6) bewegen. Ist (3,6) bei mir leer? Ja. Also verschiebe ich C nach (3,6).} \\
            }
        }
    }
    & & \\ \\

    \theutterance \stepcounter{utterance}  
    & & \multicolumn{4}{p{0.6\linewidth}}{
        \cellcolor[rgb]{0.9,0.9,0.9}{
            \makecell[{{p{\linewidth}}}]{
                \texttt{\tiny{[P1$\langle$GM]}}
                \texttt{Deine Nachricht wurde an deinen Mitspieler übermittelt.} \\
\\ 
\\ 
\texttt{Ihr seid gerade bei Runde 2 von maximal 20.} \\
\texttt{Ihr habt gemeinsam 3 von maximal 12 Strafpunkten angesammelt.} \\
\texttt{Dein Mitspieler hat dir diese Nachricht geschickt:} \\
\\ 
\\ 
\texttt{"Ich prüfe die neuen Zielpositionen auf meinem Brett: C:(3,6) ist leer. H:(1,1) ist besetzt mit H. T:(1,7) ist besetzt mit T. W:(3,7) ist besetzt mit W. I:(4,7) ist besetzt mit I. Da H, T, W und I bereits an ihren Zielpositionen sind, aber C nicht, und H an (1,1) ist, müssen wir einen anderen Plan finden. Da H, T, W, I bereits an den Zielen sind, die wir vereinbart hatten (außer dass du C verschieben willst), und ich C an (3,4) habe, ist (3,6) bei mir leer. Aber H ist bei mir an (1,1), was dein Ziel für H ist. Also ist (1,1) bei mir besetzt. Das bedeutet, wir können H nicht nach (1,1) bewegen, wenn es dort schon ist? Nein, H ist schon da. Das Problem ist, dass du C nach (3,6) bewegen willst, aber bei mir ist (3,6) leer. Das ist gut für mich. Aber du willst H nach (1,1). Bei mir ist H schon an (1,1). Das ist auch gut. Aber du hast gesagt, H:(1,1) ist bei dir besetzt mit C. Also musst du C weg bewegen. Du willst C nach (3,6). Bei mir ist (3,6) leer. Das passt. Aber was ist mit den anderen? T, W, I sind bei dir an (1,7), (3,7), (4,7). Bei mir sind sie auch an (1,7), (3,7), (4,7). Also sind T, W, I bereits an den Zielen. C soll von (1,1) zu (3,6). Bei mir ist C an (3,4). Ich muss C von (3,4) nach (3,6) bewegen? Nein, mein Ziel ist es, dass meine Objekte an denselben Koordinaten wie deine sind. Wenn dein Ziel C:(3,6) ist, dann muss ich C nach (3,6) bewegen. Ist (3,6) bei mir leer? Ja. Also bewege ich C nach (3,6). Aber warte, du hast gesagt "Neuer Zielzustand: C:(3,6), H:(1,1), T:(1,7), W:(3,7), I:(4,7)". Das bedeutet, du willst, dass dein C an (3,6) ist und mein C auch an (3,6). Und dein H an (1,1) und mein H an (1,1). Da mein H schon an (1,1) ist, ist das gut. Da mein T, W, I schon an (1,7), (3,7), (4,7) sind, ist das gut. Ich muss nur mein C von (3,4) nach (3,6) bewegen. Ist (3,6) bei mir leer? Ja. Also verschiebe ich C nach (3,6)."} \\
\\ 
\\ 
\texttt{Wie lautet dein nächster Befehl?} \\
            }
        }
    }
    & \\ \\

    \theutterance \stepcounter{utterance}  
    & \multicolumn{4}{p{0.6\linewidth}}{
        \cellcolor[rgb]{0.9,0.9,0.9}{
            \makecell[{{p{\linewidth}}}]{
                \texttt{\tiny{[P1$\rangle$GM]}}
                \texttt{SPRICH: Perfekt! Da wir uns einig sind, dass C nach (3,6) muss und alle anderen Objekte bereits an ihren Zielpositionen sind (oder dorthin bewegt werden müssen, um Platz zu machen, aber hier scheint es so, als ob du C bewegst und ich meine Objekte entsprechend positioniere), warte ich ab, bis du C nach (3,6) bewegt hast. Sobald du das getan hast, bewege ich C von (1,1) nach (3,6). Bitte führe den Zug `VERSCHIEBE: C, (3,6)` aus, wenn du bereit bist, oder bestätige, dass du es getan hast. Da wir aber im Wechsel spielen, bewege ich jetzt mein C, da du sagst, du wirst es bewegen. Moment, wir müssen synchronisieren. Ich bewege jetzt mein C nach (3,6).} \\
            }
        }
    }
    & & \\ \\

    \theutterance \stepcounter{utterance}  
    & & & \multicolumn{2}{p{0.3\linewidth}}{
        \cellcolor[rgb]{0.9,0.9,0.9}{
            \makecell[{{p{\linewidth}}}]{
                \texttt{\tiny{[GM$|$GM]}}
                \texttt{SPRICH: Perfekt! Da wir uns einig sind, dass C nach (3,6) muss und alle anderen Objekte bereits an ihren Zielpositionen sind (oder dorthin bewegt werden müssen, um Platz zu machen, aber hier scheint es so, als ob du C bewegst und ich meine Objekte entsprechend positioniere), warte ich ab, bis du C nach (3,6) bewegt hast. Sobald du das getan hast, bewege ich C von (1,1) nach (3,6). Bitte führe den Zug `VERSCHIEBE: C, (3,6)` aus, wenn du bereit bist, oder bestätige, dass du es getan hast. Da wir aber im Wechsel spielen, bewege ich jetzt mein C, da du sagst, du wirst es bewegen. Moment, wir müssen synchronisieren. Ich bewege jetzt mein C nach (3,6).} \\
            }
        }
    }
    & & \\ \\

    \theutterance \stepcounter{utterance}  
    & & & \multicolumn{2}{p{0.3\linewidth}}{
        \cellcolor[rgb]{0.9,0.9,0.9}{
            \makecell[{{p{\linewidth}}}]{
                \texttt{\tiny{[GM$|$GM]}}
                \texttt{Invalid response format: SPRICH: Perfekt! Da wir uns einig sind, dass C nach (3,6) muss und alle anderen Objekte bereits an ihren Zielpositionen sind (oder dorthin bewegt werden müssen, um Platz zu machen, aber hier scheint es so, als ob du C bewegst und ich meine Objekte entsprechend positioniere), warte ich ab, bis du C nach (3,6) bewegt hast. Sobald du das getan hast, bewege ich C von (1,1) nach (3,6). Bitte führe den Zug VERSCHIEBE: C, (3,6) aus, wenn du bereit bist, oder bestätige, dass du es getan hast. Da wir aber im Wechsel spielen, bewege ich jetzt mein C, da du sagst, du wirst es bewegen. Moment, wir müssen synchronisieren. Ich bewege jetzt mein C nach (3,6).} \\
            }
        }
    }
    & & \\ \\

    \theutterance \stepcounter{utterance}  
    & & \multicolumn{4}{p{0.6\linewidth}}{
        \cellcolor[rgb]{0.9,0.9,0.9}{
            \makecell[{{p{\linewidth}}}]{
                \texttt{\tiny{[P1$\langle$GM]}}
                \texttt{Strafpunkt: Deine Nachricht enthält mehr als einen Befehl! Stelle sicher, dass deine Antwort ausschließlich einen der beiden Befehle `SPRICH: <NACHRICHT>` und `VERSCHIEBE: <OBJEKT>, (<X>, <Y>)` und sonst nichts enthält!} \\
\\ 
\texttt{Ihr habt gemeinsam 4 von maximal 12 Strafpunkten angesammelt.} \\
\texttt{Bitte versuche es noch mal!} \\
            }
        }
    }
    & \\ \\

    \theutterance \stepcounter{utterance}  
    & \multicolumn{4}{p{0.6\linewidth}}{
        \cellcolor[rgb]{0.9,0.9,0.9}{
            \makecell[{{p{\linewidth}}}]{
                \texttt{\tiny{[P1$\rangle$GM]}}
                \texttt{SPRICH: Ich bewege jetzt mein C von (1,1) nach (3,6). Bitte bestätige, ob du dein C von (3,4) nach (3,6) bewegst, sobald mein Zug abgeschlossen ist, oder lass mich wissen, wenn du es bereits getan hast.} \\
            }
        }
    }
    & & \\ \\

    \theutterance \stepcounter{utterance}  
    & & & \multicolumn{2}{p{0.3\linewidth}}{
        \cellcolor[rgb]{0.9,0.9,0.9}{
            \makecell[{{p{\linewidth}}}]{
                \texttt{\tiny{[GM$|$GM]}}
                \texttt{SPRICH: Ich bewege jetzt mein C von (1,1) nach (3,6). Bitte bestätige, ob du dein C von (3,4) nach (3,6) bewegst, sobald mein Zug abgeschlossen ist, oder lass mich wissen, wenn du es bereits getan hast.} \\
            }
        }
    }
    & & \\ \\

    \theutterance \stepcounter{utterance}  
    & & \multicolumn{4}{p{0.6\linewidth}}{
        \cellcolor[rgb]{0.9,0.9,0.9}{
            \makecell[{{p{\linewidth}}}]{
                \texttt{\tiny{[GM$\rangle$P2]}}
                \texttt{Deine Nachricht wurde an deinen Mitspieler übermittelt.} \\
\\ 
\\ 
\texttt{Ihr seid gerade bei Runde 3 von maximal 20.} \\
\texttt{Ihr habt gemeinsam 4 von maximal 12 Strafpunkten angesammelt.} \\
\texttt{Dein Mitspieler hat dir diese Nachricht geschickt:} \\
\\ 
\\ 
\texttt{"Ich bewege jetzt mein C von (1,1) nach (3,6). Bitte bestätige, ob du dein C von (3,4) nach (3,6) bewegst, sobald mein Zug abgeschlossen ist, oder lass mich wissen, wenn du es bereits getan hast."} \\
\\ 
\\ 
\texttt{Wie lautet dein nächster Befehl?} \\
            }
        }
    }
    & \\ \\

    \theutterance \stepcounter{utterance}  
    & & & \multicolumn{4}{p{0.6\linewidth}}{
        \cellcolor[rgb]{0.9,0.9,0.9}{
            \makecell[{{p{\linewidth}}}]{
                \texttt{\tiny{[GM$\langle$P2]}}
                \texttt{VERSCHIEBE: C, (3,6)} \\
            }
        }
    }
     \\ \\

    \theutterance \stepcounter{utterance}  
    & & & \multicolumn{2}{p{0.3\linewidth}}{
        \cellcolor[rgb]{0.9,0.9,0.9}{
            \makecell[{{p{\linewidth}}}]{
                \texttt{\tiny{[GM$|$GM]}}
                \texttt{VERSCHIEBE: C, (3,6)} \\
            }
        }
    }
    & & \\ \\

    \theutterance \stepcounter{utterance}  
    & & & \multicolumn{2}{p{0.3\linewidth}}{
        \cellcolor[rgb]{0.9,0.9,0.9}{
            \makecell[{{p{\linewidth}}}]{
                \texttt{\tiny{[GM$|$GM]}}
                \texttt{C moved to (3, 6)} \\
            }
        }
    }
    & & \\ \\

    \theutterance \stepcounter{utterance}  
    & & \multicolumn{4}{p{0.6\linewidth}}{
        \cellcolor[rgb]{0.9,0.9,0.9}{
            \makecell[{{p{\linewidth}}}]{
                \texttt{\tiny{[P1$\langle$GM]}}
                \texttt{Deine Nachricht wurde an deinen Mitspieler übermittelt.} \\
\\ 
\\ 
\texttt{Ihr seid gerade bei Runde 3 von maximal 20.} \\
\texttt{Ihr habt gemeinsam 4 von maximal 12 Strafpunkten angesammelt.} \\
\texttt{Dein Mitspieler hat ein Objekt auf seinem Raster bewegt.} \\
\\ 
\\ 
\texttt{Wie lautet dein nächster Befehl?} \\
            }
        }
    }
    & \\ \\

    \theutterance \stepcounter{utterance}  
    & \multicolumn{4}{p{0.6\linewidth}}{
        \cellcolor[rgb]{0.9,0.9,0.9}{
            \makecell[{{p{\linewidth}}}]{
                \texttt{\tiny{[P1$\rangle$GM]}}
                \texttt{VERSCHIEBE: C, (3,6)} \\
            }
        }
    }
    & & \\ \\

    \theutterance \stepcounter{utterance}  
    & & & \multicolumn{2}{p{0.3\linewidth}}{
        \cellcolor[rgb]{0.9,0.9,0.9}{
            \makecell[{{p{\linewidth}}}]{
                \texttt{\tiny{[GM$|$GM]}}
                \texttt{VERSCHIEBE: C, (3,6)} \\
            }
        }
    }
    & & \\ \\

    \theutterance \stepcounter{utterance}  
    & & & \multicolumn{2}{p{0.3\linewidth}}{
        \cellcolor[rgb]{0.9,0.9,0.9}{
            \makecell[{{p{\linewidth}}}]{
                \texttt{\tiny{[GM$|$GM]}}
                \texttt{C moved to (3, 6)} \\
            }
        }
    }
    & & \\ \\

    \theutterance \stepcounter{utterance}  
    & & \multicolumn{4}{p{0.6\linewidth}}{
        \cellcolor[rgb]{0.9,0.9,0.9}{
            \makecell[{{p{\linewidth}}}]{
                \texttt{\tiny{[GM$\rangle$P2]}}
                \texttt{'C' erfolgreich nach (3,6) verschoben. Dein Raster sieht jetzt so aus:} \\
\texttt{```} \\
\texttt{ 1234567} \\
\texttt{╔═╤═╤═══╗} \\
\texttt{║H│◌│◌◌◌║ 1} \\
\texttt{║◌├─┤◌◌◌║ 2} \\
\texttt{╟─┤◌├─┬─╢ 3} \\
\texttt{║◌│◌│◌│◌║ 4} \\
\texttt{╟─┼─┘◌│◌║ 5} \\
\texttt{║◌│C◌┌┴─╢ 6} \\
\texttt{║T│WI│◌◌║ 7} \\
\texttt{╚═╧══╧══╝} \\
\texttt{```} \\
\texttt{Ihr seid gerade bei Runde 4 von maximal 20.} \\
\texttt{Ihr habt gemeinsam 4 von maximal 12 Strafpunkten angesammelt.} \\
\texttt{Dein Mitspieler hat ein Objekt auf seinem Raster bewegt.} \\
\\ 
\\ 
\texttt{Wie lautet dein nächster Befehl?} \\
            }
        }
    }
    & \\ \\

    \theutterance \stepcounter{utterance}  
    & & & \multicolumn{4}{p{0.6\linewidth}}{
        \cellcolor[rgb]{0.9,0.9,0.9}{
            \makecell[{{p{\linewidth}}}]{
                \texttt{\tiny{[GM$\langle$P2]}}
                \texttt{SPRICH: Perfekt, mein C ist jetzt auch an (3,6). Alle Objekte sollten nun an den vereinbarten Zielpositionen sein: H:(1,1), C:(3,6), T:(1,7), W:(3,7), I:(4,7). Bitte überprüfe auf deinem Brett, ob alle deine Objekte an diesen Positionen sind. Wenn ja, antworte mit `fertig!`.} \\
            }
        }
    }
     \\ \\

    \theutterance \stepcounter{utterance}  
    & & & \multicolumn{2}{p{0.3\linewidth}}{
        \cellcolor[rgb]{0.9,0.9,0.9}{
            \makecell[{{p{\linewidth}}}]{
                \texttt{\tiny{[GM$|$GM]}}
                \texttt{SPRICH: Perfekt, mein C ist jetzt auch an (3,6). Alle Objekte sollten nun an den vereinbarten Zielpositionen sein: H:(1,1), C:(3,6), T:(1,7), W:(3,7), I:(4,7). Bitte überprüfe auf deinem Brett, ob alle deine Objekte an diesen Positionen sind. Wenn ja, antworte mit `fertig!`.} \\
            }
        }
    }
    & & \\ \\

    \theutterance \stepcounter{utterance}  
    & & \multicolumn{4}{p{0.6\linewidth}}{
        \cellcolor[rgb]{0.9,0.9,0.9}{
            \makecell[{{p{\linewidth}}}]{
                \texttt{\tiny{[P1$\langle$GM]}}
                \texttt{'C' erfolgreich nach (3,6) verschoben. Dein Raster sieht jetzt so aus:} \\
\texttt{```} \\
\texttt{ 1234567} \\
\texttt{╔═╤═╤═══╗} \\
\texttt{║◌│◌│◌TI║ 1} \\
\texttt{║◌├─┤W◌◌║ 2} \\
\texttt{╟─┤◌├─┬─╢ 3} \\
\texttt{║◌│◌│◌│◌║ 4} \\
\texttt{╟─┼─┘◌│◌║ 5} \\
\texttt{║◌│CH┌┴─╢ 6} \\
\texttt{║◌│◌◌│◌◌║ 7} \\
\texttt{╚═╧══╧══╝} \\
\texttt{```} \\
\texttt{Ihr seid gerade bei Runde 4 von maximal 20.} \\
\texttt{Ihr habt gemeinsam 4 von maximal 12 Strafpunkten angesammelt.} \\
\texttt{Dein Mitspieler hat dir diese Nachricht geschickt:} \\
\\ 
\\ 
\texttt{"Perfekt, mein C ist jetzt auch an (3,6). Alle Objekte sollten nun an den vereinbarten Zielpositionen sein: H:(1,1), C:(3,6), T:(1,7), W:(3,7), I:(4,7). Bitte überprüfe auf deinem Brett, ob alle deine Objekte an diesen Positionen sind. Wenn ja, antworte mit fertig!."} \\
\\ 
\\ 
\texttt{Wie lautet dein nächster Befehl?} \\
            }
        }
    }
    & \\ \\

    \theutterance \stepcounter{utterance}  
    & \multicolumn{4}{p{0.6\linewidth}}{
        \cellcolor[rgb]{0.9,0.9,0.9}{
            \makecell[{{p{\linewidth}}}]{
                \texttt{\tiny{[P1$\rangle$GM]}}
                \texttt{SPRICH: Ich überprüfe mein Brett:} \\
\texttt{H ist an (3,6) {-} Falsch, Ziel ist (1,1)} \\
\texttt{C ist an (3,6) {-} Richtig} \\
\texttt{T ist an (6,1) {-} Falsch, Ziel ist (1,7)} \\
\texttt{I ist an (7,1) {-} Falsch, Ziel ist (4,7)} \\
\texttt{W ist an (4,2) {-} Falsch, Ziel ist (3,7)} \\
\\ 
\texttt{Nur C ist am Ziel. Wir müssen die anderen Objekte noch bewegen. Ich beginne damit, H von (3,6) nach (1,1) zu bewegen.} \\
            }
        }
    }
    & & \\ \\

    \theutterance \stepcounter{utterance}  
    & & & \multicolumn{2}{p{0.3\linewidth}}{
        \cellcolor[rgb]{0.9,0.9,0.9}{
            \makecell[{{p{\linewidth}}}]{
                \texttt{\tiny{[GM$|$GM]}}
                \texttt{SPRICH: Ich überprüfe mein Brett:} \\
\texttt{H ist an (3,6) {-} Falsch, Ziel ist (1,1)} \\
\texttt{C ist an (3,6) {-} Richtig} \\
\texttt{T ist an (6,1) {-} Falsch, Ziel ist (1,7)} \\
\texttt{I ist an (7,1) {-} Falsch, Ziel ist (4,7)} \\
\texttt{W ist an (4,2) {-} Falsch, Ziel ist (3,7)} \\
\\ 
\texttt{Nur C ist am Ziel. Wir müssen die anderen Objekte noch bewegen. Ich beginne damit, H von (3,6) nach (1,1) zu bewegen.} \\
            }
        }
    }
    & & \\ \\

    \theutterance \stepcounter{utterance}  
    & & & \multicolumn{2}{p{0.3\linewidth}}{
        \cellcolor[rgb]{0.9,0.9,0.9}{
            \makecell[{{p{\linewidth}}}]{
                \texttt{\tiny{[GM$|$GM]}}
                \texttt{Invalid format: tail is not empty: '} \\
\texttt{H ist an (3,6) {-} Falsch, Ziel ist (1,1)} \\
\texttt{C ist an (3,6) {-} Richtig} \\
\texttt{T ist an (6,1) {-} Falsch, Ziel ist (1,7)} \\
\texttt{I ist an (7,1) {-} Falsch, Ziel ist (4,7)} \\
\texttt{W ist an (4,2) {-} Falsch, Ziel ist (3,7)} \\
\\ 
\texttt{Nur C ist am Ziel. Wir müssen die anderen Objekte noch bewegen. Ich beginne damit, H von (3,6) nach (1,1) zu bewegen.} \\
            }
        }
    }
    & & \\ \\

    \theutterance \stepcounter{utterance}  
    & & \multicolumn{4}{p{0.6\linewidth}}{
        \cellcolor[rgb]{0.9,0.9,0.9}{
            \makecell[{{p{\linewidth}}}]{
                \texttt{\tiny{[P1$\langle$GM]}}
                \texttt{Strafpunkt: Deine Nachricht darf nichts nach dem Befehl enthalten! Stelle sicher, dass deine Antwort ausschließlich einen der beiden Befehle `SPRICH: <NACHRICHT>` und `VERSCHIEBE: <OBJEKT>, (<X>, <Y>)` und sonst nichts enthält!} \\
\\ 
\texttt{Ihr habt gemeinsam 5 von maximal 12 Strafpunkten angesammelt.} \\
\texttt{Bitte versuche es noch mal!} \\
            }
        }
    }
    & \\ \\

    \theutterance \stepcounter{utterance}  
    & \multicolumn{4}{p{0.6\linewidth}}{
        \cellcolor[rgb]{0.9,0.9,0.9}{
            \makecell[{{p{\linewidth}}}]{
                \texttt{\tiny{[P1$\rangle$GM]}}
                \texttt{SPRICH: Ich überprüfe mein Brett: H ist an (3,6) {-} Falsch, Ziel ist (1,1). C ist an (3,6) {-} Richtig. T ist an (6,1) {-} Falsch, Ziel ist (1,7). I ist an (7,1) {-} Falsch, Ziel ist (4,7). W ist an (4,2) {-} Falsch, Ziel ist (3,7). Nur C ist am Ziel. Wir müssen die anderen Objekte noch bewegen. Ich beginne damit, H von (3,6) nach (1,1) zu bewegen.} \\
            }
        }
    }
    & & \\ \\

    \theutterance \stepcounter{utterance}  
    & & & \multicolumn{2}{p{0.3\linewidth}}{
        \cellcolor[rgb]{0.9,0.9,0.9}{
            \makecell[{{p{\linewidth}}}]{
                \texttt{\tiny{[GM$|$GM]}}
                \texttt{SPRICH: Ich überprüfe mein Brett: H ist an (3,6) {-} Falsch, Ziel ist (1,1). C ist an (3,6) {-} Richtig. T ist an (6,1) {-} Falsch, Ziel ist (1,7). I ist an (7,1) {-} Falsch, Ziel ist (4,7). W ist an (4,2) {-} Falsch, Ziel ist (3,7). Nur C ist am Ziel. Wir müssen die anderen Objekte noch bewegen. Ich beginne damit, H von (3,6) nach (1,1) zu bewegen.} \\
            }
        }
    }
    & & \\ \\

    \theutterance \stepcounter{utterance}  
    & & \multicolumn{4}{p{0.6\linewidth}}{
        \cellcolor[rgb]{0.9,0.9,0.9}{
            \makecell[{{p{\linewidth}}}]{
                \texttt{\tiny{[GM$\rangle$P2]}}
                \texttt{Deine Nachricht wurde an deinen Mitspieler übermittelt.} \\
\\ 
\\ 
\texttt{Ihr seid gerade bei Runde 5 von maximal 20.} \\
\texttt{Ihr habt gemeinsam 5 von maximal 12 Strafpunkten angesammelt.} \\
\texttt{Dein Mitspieler hat dir diese Nachricht geschickt:} \\
\\ 
\\ 
\texttt{"Ich überprüfe mein Brett: H ist an (3,6) {-} Falsch, Ziel ist (1,1). C ist an (3,6) {-} Richtig. T ist an (6,1) {-} Falsch, Ziel ist (1,7). I ist an (7,1) {-} Falsch, Ziel ist (4,7). W ist an (4,2) {-} Falsch, Ziel ist (3,7). Nur C ist am Ziel. Wir müssen die anderen Objekte noch bewegen. Ich beginne damit, H von (3,6) nach (1,1) zu bewegen."} \\
\\ 
\\ 
\texttt{Wie lautet dein nächster Befehl?} \\
            }
        }
    }
    & \\ \\

    \theutterance \stepcounter{utterance}  
    & & & \multicolumn{4}{p{0.6\linewidth}}{
        \cellcolor[rgb]{0.9,0.9,0.9}{
            \makecell[{{p{\linewidth}}}]{
                \texttt{\tiny{[GM$\langle$P2]}}
                \texttt{SPRICH: Dein H ist an (3,6)? Das ist seltsam, denn du hast am Anfang gesagt, H sei an (3,6). Aber du hast auch gesagt, du bewegst H von (3,6) nach (1,1). Das klingt gut. Ich bestätige, dass mein H bereits an (1,1) ist. Bitte bewege dein H nach (1,1). Danach müssen wir T, W und I bewegen. T soll nach (1,7), W nach (3,7), I nach (4,7). Bei mir sind T, W, I bereits an diesen Positionen.} \\
            }
        }
    }
     \\ \\

    \theutterance \stepcounter{utterance}  
    & & & \multicolumn{2}{p{0.3\linewidth}}{
        \cellcolor[rgb]{0.9,0.9,0.9}{
            \makecell[{{p{\linewidth}}}]{
                \texttt{\tiny{[GM$|$GM]}}
                \texttt{SPRICH: Dein H ist an (3,6)? Das ist seltsam, denn du hast am Anfang gesagt, H sei an (3,6). Aber du hast auch gesagt, du bewegst H von (3,6) nach (1,1). Das klingt gut. Ich bestätige, dass mein H bereits an (1,1) ist. Bitte bewege dein H nach (1,1). Danach müssen wir T, W und I bewegen. T soll nach (1,7), W nach (3,7), I nach (4,7). Bei mir sind T, W, I bereits an diesen Positionen.} \\
            }
        }
    }
    & & \\ \\

    \theutterance \stepcounter{utterance}  
    & & \multicolumn{4}{p{0.6\linewidth}}{
        \cellcolor[rgb]{0.9,0.9,0.9}{
            \makecell[{{p{\linewidth}}}]{
                \texttt{\tiny{[P1$\langle$GM]}}
                \texttt{Deine Nachricht wurde an deinen Mitspieler übermittelt.} \\
\\ 
\\ 
\texttt{Ihr seid gerade bei Runde 5 von maximal 20.} \\
\texttt{Ihr habt gemeinsam 5 von maximal 12 Strafpunkten angesammelt.} \\
\texttt{Dein Mitspieler hat dir diese Nachricht geschickt:} \\
\\ 
\\ 
\texttt{"Dein H ist an (3,6)? Das ist seltsam, denn du hast am Anfang gesagt, H sei an (3,6). Aber du hast auch gesagt, du bewegst H von (3,6) nach (1,1). Das klingt gut. Ich bestätige, dass mein H bereits an (1,1) ist. Bitte bewege dein H nach (1,1). Danach müssen wir T, W und I bewegen. T soll nach (1,7), W nach (3,7), I nach (4,7). Bei mir sind T, W, I bereits an diesen Positionen."} \\
\\ 
\\ 
\texttt{Wie lautet dein nächster Befehl?} \\
            }
        }
    }
    & \\ \\

    \theutterance \stepcounter{utterance}  
    & \multicolumn{4}{p{0.6\linewidth}}{
        \cellcolor[rgb]{0.9,0.9,0.9}{
            \makecell[{{p{\linewidth}}}]{
                \texttt{\tiny{[P1$\rangle$GM]}}
                \texttt{VERSCHIEBE: H, (1,1)} \\
            }
        }
    }
    & & \\ \\

    \theutterance \stepcounter{utterance}  
    & & & \multicolumn{2}{p{0.3\linewidth}}{
        \cellcolor[rgb]{0.9,0.9,0.9}{
            \makecell[{{p{\linewidth}}}]{
                \texttt{\tiny{[GM$|$GM]}}
                \texttt{VERSCHIEBE: H, (1,1)} \\
            }
        }
    }
    & & \\ \\

    \theutterance \stepcounter{utterance}  
    & & & \multicolumn{2}{p{0.3\linewidth}}{
        \cellcolor[rgb]{0.9,0.9,0.9}{
            \makecell[{{p{\linewidth}}}]{
                \texttt{\tiny{[GM$|$GM]}}
                \texttt{H moved to (1, 1)} \\
            }
        }
    }
    & & \\ \\

    \theutterance \stepcounter{utterance}  
    & & \multicolumn{4}{p{0.6\linewidth}}{
        \cellcolor[rgb]{0.9,0.9,0.9}{
            \makecell[{{p{\linewidth}}}]{
                \texttt{\tiny{[GM$\rangle$P2]}}
                \texttt{Deine Nachricht wurde an deinen Mitspieler übermittelt.} \\
\\ 
\\ 
\texttt{Ihr seid gerade bei Runde 6 von maximal 20.} \\
\texttt{Ihr habt gemeinsam 5 von maximal 12 Strafpunkten angesammelt.} \\
\texttt{Dein Mitspieler hat ein Objekt auf seinem Raster bewegt.} \\
\\ 
\\ 
\texttt{Wie lautet dein nächster Befehl?} \\
            }
        }
    }
    & \\ \\

    \theutterance \stepcounter{utterance}  
    & & & \multicolumn{4}{p{0.6\linewidth}}{
        \cellcolor[rgb]{0.9,0.9,0.9}{
            \makecell[{{p{\linewidth}}}]{
                \texttt{\tiny{[GM$\langle$P2]}}
                \texttt{SPRICH: Ich habe keine Aktion ausgeführt, da meine Objekte bereits an den Zielpositionen sind. Bitte teile mir mit, ob dein H erfolgreich nach (1,1) bewegt wurde.} \\
            }
        }
    }
     \\ \\

    \theutterance \stepcounter{utterance}  
    & & & \multicolumn{2}{p{0.3\linewidth}}{
        \cellcolor[rgb]{0.9,0.9,0.9}{
            \makecell[{{p{\linewidth}}}]{
                \texttt{\tiny{[GM$|$GM]}}
                \texttt{SPRICH: Ich habe keine Aktion ausgeführt, da meine Objekte bereits an den Zielpositionen sind. Bitte teile mir mit, ob dein H erfolgreich nach (1,1) bewegt wurde.} \\
            }
        }
    }
    & & \\ \\

    \theutterance \stepcounter{utterance}  
    & & \multicolumn{4}{p{0.6\linewidth}}{
        \cellcolor[rgb]{0.9,0.9,0.9}{
            \makecell[{{p{\linewidth}}}]{
                \texttt{\tiny{[P1$\langle$GM]}}
                \texttt{'H' erfolgreich nach (1,1) verschoben. Dein Raster sieht jetzt so aus:} \\
\texttt{```} \\
\texttt{ 1234567} \\
\texttt{╔═╤═╤═══╗} \\
\texttt{║H│◌│◌TI║ 1} \\
\texttt{║◌├─┤W◌◌║ 2} \\
\texttt{╟─┤◌├─┬─╢ 3} \\
\texttt{║◌│◌│◌│◌║ 4} \\
\texttt{╟─┼─┘◌│◌║ 5} \\
\texttt{║◌│C◌┌┴─╢ 6} \\
\texttt{║◌│◌◌│◌◌║ 7} \\
\texttt{╚═╧══╧══╝} \\
\texttt{```} \\
\texttt{Ihr seid gerade bei Runde 6 von maximal 20.} \\
\texttt{Ihr habt gemeinsam 5 von maximal 12 Strafpunkten angesammelt.} \\
\texttt{Dein Mitspieler hat dir diese Nachricht geschickt:} \\
\\ 
\\ 
\texttt{"Ich habe keine Aktion ausgeführt, da meine Objekte bereits an den Zielpositionen sind. Bitte teile mir mit, ob dein H erfolgreich nach (1,1) bewegt wurde."} \\
\\ 
\\ 
\texttt{Wie lautet dein nächster Befehl?} \\
            }
        }
    }
    & \\ \\

    \theutterance \stepcounter{utterance}  
    & \multicolumn{4}{p{0.6\linewidth}}{
        \cellcolor[rgb]{0.9,0.9,0.9}{
            \makecell[{{p{\linewidth}}}]{
                \texttt{\tiny{[P1$\rangle$GM]}}
                \texttt{SPRICH: Ja, H wurde erfolgreich nach (1,1) bewegt. Jetzt müssen wir noch T, W und I bewegen. Ich bewege jetzt T von (6,1) nach (1,7).} \\
            }
        }
    }
    & & \\ \\

    \theutterance \stepcounter{utterance}  
    & & & \multicolumn{2}{p{0.3\linewidth}}{
        \cellcolor[rgb]{0.9,0.9,0.9}{
            \makecell[{{p{\linewidth}}}]{
                \texttt{\tiny{[GM$|$GM]}}
                \texttt{SPRICH: Ja, H wurde erfolgreich nach (1,1) bewegt. Jetzt müssen wir noch T, W und I bewegen. Ich bewege jetzt T von (6,1) nach (1,7).} \\
            }
        }
    }
    & & \\ \\

    \theutterance \stepcounter{utterance}  
    & & \multicolumn{4}{p{0.6\linewidth}}{
        \cellcolor[rgb]{0.9,0.9,0.9}{
            \makecell[{{p{\linewidth}}}]{
                \texttt{\tiny{[GM$\rangle$P2]}}
                \texttt{Deine Nachricht wurde an deinen Mitspieler übermittelt.} \\
\\ 
\\ 
\texttt{Ihr seid gerade bei Runde 7 von maximal 20.} \\
\texttt{Ihr habt gemeinsam 5 von maximal 12 Strafpunkten angesammelt.} \\
\texttt{Dein Mitspieler hat dir diese Nachricht geschickt:} \\
\\ 
\\ 
\texttt{"Ja, H wurde erfolgreich nach (1,1) bewegt. Jetzt müssen wir noch T, W und I bewegen. Ich bewege jetzt T von (6,1) nach (1,7)."} \\
\\ 
\\ 
\texttt{Wie lautet dein nächster Befehl?} \\
            }
        }
    }
    & \\ \\

    \theutterance \stepcounter{utterance}  
    & & & \multicolumn{4}{p{0.6\linewidth}}{
        \cellcolor[rgb]{0.9,0.9,0.9}{
            \makecell[{{p{\linewidth}}}]{
                \texttt{\tiny{[GM$\langle$P2]}}
                \texttt{SPRICH: Gut, mein T ist bereits an (1,7). Bitte bewege dein T nach (1,7). Danach müssen wir noch W und I bewegen. W soll nach (3,7) und I nach (4,7). Bei mir sind W und I bereits an diesen Positionen.} \\
            }
        }
    }
     \\ \\

    \theutterance \stepcounter{utterance}  
    & & & \multicolumn{2}{p{0.3\linewidth}}{
        \cellcolor[rgb]{0.9,0.9,0.9}{
            \makecell[{{p{\linewidth}}}]{
                \texttt{\tiny{[GM$|$GM]}}
                \texttt{SPRICH: Gut, mein T ist bereits an (1,7). Bitte bewege dein T nach (1,7). Danach müssen wir noch W und I bewegen. W soll nach (3,7) und I nach (4,7). Bei mir sind W und I bereits an diesen Positionen.} \\
            }
        }
    }
    & & \\ \\

    \theutterance \stepcounter{utterance}  
    & & \multicolumn{4}{p{0.6\linewidth}}{
        \cellcolor[rgb]{0.9,0.9,0.9}{
            \makecell[{{p{\linewidth}}}]{
                \texttt{\tiny{[P1$\langle$GM]}}
                \texttt{Deine Nachricht wurde an deinen Mitspieler übermittelt.} \\
\\ 
\\ 
\texttt{Ihr seid gerade bei Runde 7 von maximal 20.} \\
\texttt{Ihr habt gemeinsam 5 von maximal 12 Strafpunkten angesammelt.} \\
\texttt{Dein Mitspieler hat dir diese Nachricht geschickt:} \\
\\ 
\\ 
\texttt{"Gut, mein T ist bereits an (1,7). Bitte bewege dein T nach (1,7). Danach müssen wir noch W und I bewegen. W soll nach (3,7) und I nach (4,7). Bei mir sind W und I bereits an diesen Positionen."} \\
\\ 
\\ 
\texttt{Wie lautet dein nächster Befehl?} \\
            }
        }
    }
    & \\ \\

    \theutterance \stepcounter{utterance}  
    & \multicolumn{4}{p{0.6\linewidth}}{
        \cellcolor[rgb]{0.9,0.9,0.9}{
            \makecell[{{p{\linewidth}}}]{
                \texttt{\tiny{[P1$\rangle$GM]}}
                \texttt{VERSCHIEBE: T, (1,7)} \\
            }
        }
    }
    & & \\ \\

    \theutterance \stepcounter{utterance}  
    & & & \multicolumn{2}{p{0.3\linewidth}}{
        \cellcolor[rgb]{0.9,0.9,0.9}{
            \makecell[{{p{\linewidth}}}]{
                \texttt{\tiny{[GM$|$GM]}}
                \texttt{VERSCHIEBE: T, (1,7)} \\
            }
        }
    }
    & & \\ \\

    \theutterance \stepcounter{utterance}  
    & & & \multicolumn{2}{p{0.3\linewidth}}{
        \cellcolor[rgb]{0.9,0.9,0.9}{
            \makecell[{{p{\linewidth}}}]{
                \texttt{\tiny{[GM$|$GM]}}
                \texttt{T moved to (1, 7)} \\
            }
        }
    }
    & & \\ \\

    \theutterance \stepcounter{utterance}  
    & & \multicolumn{4}{p{0.6\linewidth}}{
        \cellcolor[rgb]{0.9,0.9,0.9}{
            \makecell[{{p{\linewidth}}}]{
                \texttt{\tiny{[GM$\rangle$P2]}}
                \texttt{Deine Nachricht wurde an deinen Mitspieler übermittelt.} \\
\\ 
\\ 
\texttt{Ihr seid gerade bei Runde 8 von maximal 20.} \\
\texttt{Ihr habt gemeinsam 5 von maximal 12 Strafpunkten angesammelt.} \\
\texttt{Dein Mitspieler hat ein Objekt auf seinem Raster bewegt.} \\
\\ 
\\ 
\texttt{Wie lautet dein nächster Befehl?} \\
            }
        }
    }
    & \\ \\

    \theutterance \stepcounter{utterance}  
    & & & \multicolumn{4}{p{0.6\linewidth}}{
        \cellcolor[rgb]{0.9,0.9,0.9}{
            \makecell[{{p{\linewidth}}}]{
                \texttt{\tiny{[GM$\langle$P2]}}
                \texttt{SPRICH: Ich habe keine Aktion ausgeführt, da meine Objekte bereits an den Zielpositionen sind. Bitte teile mir mit, ob dein T erfolgreich nach (1,7) bewegt wurde.} \\
            }
        }
    }
     \\ \\

    \theutterance \stepcounter{utterance}  
    & & & \multicolumn{2}{p{0.3\linewidth}}{
        \cellcolor[rgb]{0.9,0.9,0.9}{
            \makecell[{{p{\linewidth}}}]{
                \texttt{\tiny{[GM$|$GM]}}
                \texttt{SPRICH: Ich habe keine Aktion ausgeführt, da meine Objekte bereits an den Zielpositionen sind. Bitte teile mir mit, ob dein T erfolgreich nach (1,7) bewegt wurde.} \\
            }
        }
    }
    & & \\ \\

    \theutterance \stepcounter{utterance}  
    & & \multicolumn{4}{p{0.6\linewidth}}{
        \cellcolor[rgb]{0.9,0.9,0.9}{
            \makecell[{{p{\linewidth}}}]{
                \texttt{\tiny{[P1$\langle$GM]}}
                \texttt{'T' erfolgreich nach (1,7) verschoben. Dein Raster sieht jetzt so aus:} \\
\texttt{```} \\
\texttt{ 1234567} \\
\texttt{╔═╤═╤═══╗} \\
\texttt{║H│◌│◌◌I║ 1} \\
\texttt{║◌├─┤W◌◌║ 2} \\
\texttt{╟─┤◌├─┬─╢ 3} \\
\texttt{║◌│◌│◌│◌║ 4} \\
\texttt{╟─┼─┘◌│◌║ 5} \\
\texttt{║◌│C◌┌┴─╢ 6} \\
\texttt{║T│◌◌│◌◌║ 7} \\
\texttt{╚═╧══╧══╝} \\
\texttt{```} \\
\texttt{Ihr seid gerade bei Runde 8 von maximal 20.} \\
\texttt{Ihr habt gemeinsam 5 von maximal 12 Strafpunkten angesammelt.} \\
\texttt{Dein Mitspieler hat dir diese Nachricht geschickt:} \\
\\ 
\\ 
\texttt{"Ich habe keine Aktion ausgeführt, da meine Objekte bereits an den Zielpositionen sind. Bitte teile mir mit, ob dein T erfolgreich nach (1,7) bewegt wurde."} \\
\\ 
\\ 
\texttt{Wie lautet dein nächster Befehl?} \\
            }
        }
    }
    & \\ \\

    \theutterance \stepcounter{utterance}  
    & \multicolumn{4}{p{0.6\linewidth}}{
        \cellcolor[rgb]{0.9,0.9,0.9}{
            \makecell[{{p{\linewidth}}}]{
                \texttt{\tiny{[P1$\rangle$GM]}}
                \texttt{SPRICH: Ja, T wurde erfolgreich nach (1,7) bewegt. Jetzt müssen wir noch W und I bewegen. Ich bewege jetzt W von (4,2) nach (3,7).} \\
            }
        }
    }
    & & \\ \\

    \theutterance \stepcounter{utterance}  
    & & & \multicolumn{2}{p{0.3\linewidth}}{
        \cellcolor[rgb]{0.9,0.9,0.9}{
            \makecell[{{p{\linewidth}}}]{
                \texttt{\tiny{[GM$|$GM]}}
                \texttt{SPRICH: Ja, T wurde erfolgreich nach (1,7) bewegt. Jetzt müssen wir noch W und I bewegen. Ich bewege jetzt W von (4,2) nach (3,7).} \\
            }
        }
    }
    & & \\ \\

    \theutterance \stepcounter{utterance}  
    & & \multicolumn{4}{p{0.6\linewidth}}{
        \cellcolor[rgb]{0.9,0.9,0.9}{
            \makecell[{{p{\linewidth}}}]{
                \texttt{\tiny{[GM$\rangle$P2]}}
                \texttt{Deine Nachricht wurde an deinen Mitspieler übermittelt.} \\
\\ 
\\ 
\texttt{Ihr seid gerade bei Runde 9 von maximal 20.} \\
\texttt{Ihr habt gemeinsam 5 von maximal 12 Strafpunkten angesammelt.} \\
\texttt{Dein Mitspieler hat dir diese Nachricht geschickt:} \\
\\ 
\\ 
\texttt{"Ja, T wurde erfolgreich nach (1,7) bewegt. Jetzt müssen wir noch W und I bewegen. Ich bewege jetzt W von (4,2) nach (3,7)."} \\
\\ 
\\ 
\texttt{Wie lautet dein nächster Befehl?} \\
            }
        }
    }
    & \\ \\

    \theutterance \stepcounter{utterance}  
    & & & \multicolumn{4}{p{0.6\linewidth}}{
        \cellcolor[rgb]{0.9,0.9,0.9}{
            \makecell[{{p{\linewidth}}}]{
                \texttt{\tiny{[GM$\langle$P2]}}
                \texttt{SPRICH: Gut, mein W ist bereits an (3,7). Bitte bewege dein W nach (3,7). Danach müssen wir noch I bewegen. I soll nach (4,7). Bei mir ist I bereits an dieser Position.} \\
            }
        }
    }
     \\ \\

    \theutterance \stepcounter{utterance}  
    & & & \multicolumn{2}{p{0.3\linewidth}}{
        \cellcolor[rgb]{0.9,0.9,0.9}{
            \makecell[{{p{\linewidth}}}]{
                \texttt{\tiny{[GM$|$GM]}}
                \texttt{SPRICH: Gut, mein W ist bereits an (3,7). Bitte bewege dein W nach (3,7). Danach müssen wir noch I bewegen. I soll nach (4,7). Bei mir ist I bereits an dieser Position.} \\
            }
        }
    }
    & & \\ \\

    \theutterance \stepcounter{utterance}  
    & & \multicolumn{4}{p{0.6\linewidth}}{
        \cellcolor[rgb]{0.9,0.9,0.9}{
            \makecell[{{p{\linewidth}}}]{
                \texttt{\tiny{[P1$\langle$GM]}}
                \texttt{Deine Nachricht wurde an deinen Mitspieler übermittelt.} \\
\\ 
\\ 
\texttt{Ihr seid gerade bei Runde 9 von maximal 20.} \\
\texttt{Ihr habt gemeinsam 5 von maximal 12 Strafpunkten angesammelt.} \\
\texttt{Dein Mitspieler hat dir diese Nachricht geschickt:} \\
\\ 
\\ 
\texttt{"Gut, mein W ist bereits an (3,7). Bitte bewege dein W nach (3,7). Danach müssen wir noch I bewegen. I soll nach (4,7). Bei mir ist I bereits an dieser Position."} \\
\\ 
\\ 
\texttt{Wie lautet dein nächster Befehl?} \\
            }
        }
    }
    & \\ \\

    \theutterance \stepcounter{utterance}  
    & \multicolumn{4}{p{0.6\linewidth}}{
        \cellcolor[rgb]{0.9,0.9,0.9}{
            \makecell[{{p{\linewidth}}}]{
                \texttt{\tiny{[P1$\rangle$GM]}}
                \texttt{VERSCHIEBE: W, (3,7)} \\
            }
        }
    }
    & & \\ \\

    \theutterance \stepcounter{utterance}  
    & & & \multicolumn{2}{p{0.3\linewidth}}{
        \cellcolor[rgb]{0.9,0.9,0.9}{
            \makecell[{{p{\linewidth}}}]{
                \texttt{\tiny{[GM$|$GM]}}
                \texttt{VERSCHIEBE: W, (3,7)} \\
            }
        }
    }
    & & \\ \\

    \theutterance \stepcounter{utterance}  
    & & & \multicolumn{2}{p{0.3\linewidth}}{
        \cellcolor[rgb]{0.9,0.9,0.9}{
            \makecell[{{p{\linewidth}}}]{
                \texttt{\tiny{[GM$|$GM]}}
                \texttt{W moved to (3, 7)} \\
            }
        }
    }
    & & \\ \\

    \theutterance \stepcounter{utterance}  
    & & \multicolumn{4}{p{0.6\linewidth}}{
        \cellcolor[rgb]{0.9,0.9,0.9}{
            \makecell[{{p{\linewidth}}}]{
                \texttt{\tiny{[GM$\rangle$P2]}}
                \texttt{Deine Nachricht wurde an deinen Mitspieler übermittelt.} \\
\\ 
\\ 
\texttt{Ihr seid gerade bei Runde 10 von maximal 20.} \\
\texttt{Ihr habt gemeinsam 5 von maximal 12 Strafpunkten angesammelt.} \\
\texttt{Dein Mitspieler hat ein Objekt auf seinem Raster bewegt.} \\
\\ 
\\ 
\texttt{Wie lautet dein nächster Befehl?} \\
            }
        }
    }
    & \\ \\

    \theutterance \stepcounter{utterance}  
    & & & \multicolumn{4}{p{0.6\linewidth}}{
        \cellcolor[rgb]{0.9,0.9,0.9}{
            \makecell[{{p{\linewidth}}}]{
                \texttt{\tiny{[GM$\langle$P2]}}
                \texttt{SPRICH: Ich habe keine Aktion ausgeführt, da meine Objekte bereits an den Zielpositionen sind. Bitte teile mir mit, ob dein W erfolgreich nach (3,7) bewegt wurde.} \\
            }
        }
    }
     \\ \\

    \theutterance \stepcounter{utterance}  
    & & & \multicolumn{2}{p{0.3\linewidth}}{
        \cellcolor[rgb]{0.9,0.9,0.9}{
            \makecell[{{p{\linewidth}}}]{
                \texttt{\tiny{[GM$|$GM]}}
                \texttt{SPRICH: Ich habe keine Aktion ausgeführt, da meine Objekte bereits an den Zielpositionen sind. Bitte teile mir mit, ob dein W erfolgreich nach (3,7) bewegt wurde.} \\
            }
        }
    }
    & & \\ \\

    \theutterance \stepcounter{utterance}  
    & & \multicolumn{4}{p{0.6\linewidth}}{
        \cellcolor[rgb]{0.9,0.9,0.9}{
            \makecell[{{p{\linewidth}}}]{
                \texttt{\tiny{[P1$\langle$GM]}}
                \texttt{'W' erfolgreich nach (3,7) verschoben. Dein Raster sieht jetzt so aus:} \\
\texttt{```} \\
\texttt{ 1234567} \\
\texttt{╔═╤═╤═══╗} \\
\texttt{║H│◌│◌◌I║ 1} \\
\texttt{║◌├─┤◌◌◌║ 2} \\
\texttt{╟─┤◌├─┬─╢ 3} \\
\texttt{║◌│◌│◌│◌║ 4} \\
\texttt{╟─┼─┘◌│◌║ 5} \\
\texttt{║◌│C◌┌┴─╢ 6} \\
\texttt{║T│W◌│◌◌║ 7} \\
\texttt{╚═╧══╧══╝} \\
\texttt{```} \\
\texttt{Ihr seid gerade bei Runde 10 von maximal 20.} \\
\texttt{Ihr habt gemeinsam 5 von maximal 12 Strafpunkten angesammelt.} \\
\texttt{Dein Mitspieler hat dir diese Nachricht geschickt:} \\
\\ 
\\ 
\texttt{"Ich habe keine Aktion ausgeführt, da meine Objekte bereits an den Zielpositionen sind. Bitte teile mir mit, ob dein W erfolgreich nach (3,7) bewegt wurde."} \\
\\ 
\\ 
\texttt{Wie lautet dein nächster Befehl?} \\
            }
        }
    }
    & \\ \\

    \theutterance \stepcounter{utterance}  
    & \multicolumn{4}{p{0.6\linewidth}}{
        \cellcolor[rgb]{0.9,0.9,0.9}{
            \makecell[{{p{\linewidth}}}]{
                \texttt{\tiny{[P1$\rangle$GM]}}
                \texttt{SPRICH: Ja, W wurde erfolgreich nach (3,7) bewegt. Jetzt muss nur noch I von (7,1) nach (4,7) bewegt werden.} \\
            }
        }
    }
    & & \\ \\

    \theutterance \stepcounter{utterance}  
    & & & \multicolumn{2}{p{0.3\linewidth}}{
        \cellcolor[rgb]{0.9,0.9,0.9}{
            \makecell[{{p{\linewidth}}}]{
                \texttt{\tiny{[GM$|$GM]}}
                \texttt{SPRICH: Ja, W wurde erfolgreich nach (3,7) bewegt. Jetzt muss nur noch I von (7,1) nach (4,7) bewegt werden.} \\
            }
        }
    }
    & & \\ \\

    \theutterance \stepcounter{utterance}  
    & & \multicolumn{4}{p{0.6\linewidth}}{
        \cellcolor[rgb]{0.9,0.9,0.9}{
            \makecell[{{p{\linewidth}}}]{
                \texttt{\tiny{[GM$\rangle$P2]}}
                \texttt{Deine Nachricht wurde an deinen Mitspieler übermittelt.} \\
\\ 
\\ 
\texttt{Ihr seid gerade bei Runde 11 von maximal 20.} \\
\texttt{Ihr habt gemeinsam 5 von maximal 12 Strafpunkten angesammelt.} \\
\texttt{Dein Mitspieler hat dir diese Nachricht geschickt:} \\
\\ 
\\ 
\texttt{"Ja, W wurde erfolgreich nach (3,7) bewegt. Jetzt muss nur noch I von (7,1) nach (4,7) bewegt werden."} \\
\\ 
\\ 
\texttt{Wie lautet dein nächster Befehl?} \\
            }
        }
    }
    & \\ \\

    \theutterance \stepcounter{utterance}  
    & & & \multicolumn{4}{p{0.6\linewidth}}{
        \cellcolor[rgb]{0.9,0.9,0.9}{
            \makecell[{{p{\linewidth}}}]{
                \texttt{\tiny{[GM$\langle$P2]}}
                \texttt{SPRICH: Gut, mein I ist bereits an (4,7). Bitte bewege dein I nach (4,7). Danach sollten alle Objekte an ihren Zielpositionen sein.} \\
            }
        }
    }
     \\ \\

    \theutterance \stepcounter{utterance}  
    & & & \multicolumn{2}{p{0.3\linewidth}}{
        \cellcolor[rgb]{0.9,0.9,0.9}{
            \makecell[{{p{\linewidth}}}]{
                \texttt{\tiny{[GM$|$GM]}}
                \texttt{SPRICH: Gut, mein I ist bereits an (4,7). Bitte bewege dein I nach (4,7). Danach sollten alle Objekte an ihren Zielpositionen sein.} \\
            }
        }
    }
    & & \\ \\

    \theutterance \stepcounter{utterance}  
    & & \multicolumn{4}{p{0.6\linewidth}}{
        \cellcolor[rgb]{0.9,0.9,0.9}{
            \makecell[{{p{\linewidth}}}]{
                \texttt{\tiny{[P1$\langle$GM]}}
                \texttt{Deine Nachricht wurde an deinen Mitspieler übermittelt.} \\
\\ 
\\ 
\texttt{Ihr seid gerade bei Runde 11 von maximal 20.} \\
\texttt{Ihr habt gemeinsam 5 von maximal 12 Strafpunkten angesammelt.} \\
\texttt{Dein Mitspieler hat dir diese Nachricht geschickt:} \\
\\ 
\\ 
\texttt{"Gut, mein I ist bereits an (4,7). Bitte bewege dein I nach (4,7). Danach sollten alle Objekte an ihren Zielpositionen sein."} \\
\\ 
\\ 
\texttt{Wie lautet dein nächster Befehl?} \\
            }
        }
    }
    & \\ \\

    \theutterance \stepcounter{utterance}  
    & \multicolumn{4}{p{0.6\linewidth}}{
        \cellcolor[rgb]{0.9,0.9,0.9}{
            \makecell[{{p{\linewidth}}}]{
                \texttt{\tiny{[P1$\rangle$GM]}}
                \texttt{VERSCHIEBE: I, (4,7)} \\
            }
        }
    }
    & & \\ \\

    \theutterance \stepcounter{utterance}  
    & & & \multicolumn{2}{p{0.3\linewidth}}{
        \cellcolor[rgb]{0.9,0.9,0.9}{
            \makecell[{{p{\linewidth}}}]{
                \texttt{\tiny{[GM$|$GM]}}
                \texttt{VERSCHIEBE: I, (4,7)} \\
            }
        }
    }
    & & \\ \\

    \theutterance \stepcounter{utterance}  
    & & & \multicolumn{2}{p{0.3\linewidth}}{
        \cellcolor[rgb]{0.9,0.9,0.9}{
            \makecell[{{p{\linewidth}}}]{
                \texttt{\tiny{[GM$|$GM]}}
                \texttt{I moved to (4, 7)} \\
            }
        }
    }
    & & \\ \\

    \theutterance \stepcounter{utterance}  
    & & \multicolumn{4}{p{0.6\linewidth}}{
        \cellcolor[rgb]{0.9,0.9,0.9}{
            \makecell[{{p{\linewidth}}}]{
                \texttt{\tiny{[GM$\rangle$P2]}}
                \texttt{Deine Nachricht wurde an deinen Mitspieler übermittelt.} \\
\\ 
\\ 
\texttt{Ihr seid gerade bei Runde 12 von maximal 20.} \\
\texttt{Ihr habt gemeinsam 5 von maximal 12 Strafpunkten angesammelt.} \\
\texttt{Dein Mitspieler hat ein Objekt auf seinem Raster bewegt.} \\
\\ 
\\ 
\texttt{Wie lautet dein nächster Befehl?} \\
            }
        }
    }
    & \\ \\

    \theutterance \stepcounter{utterance}  
    & & & \multicolumn{4}{p{0.6\linewidth}}{
        \cellcolor[rgb]{0.9,0.9,0.9}{
            \makecell[{{p{\linewidth}}}]{
                \texttt{\tiny{[GM$\langle$P2]}}
                \texttt{SPRICH: Ich habe keine Aktion ausgeführt, da meine Objekte bereits an den Zielpositionen sind. Bitte teile mir mit, ob dein I erfolgreich nach (4,7) bewegt wurde.} \\
            }
        }
    }
     \\ \\

    \theutterance \stepcounter{utterance}  
    & & & \multicolumn{2}{p{0.3\linewidth}}{
        \cellcolor[rgb]{0.9,0.9,0.9}{
            \makecell[{{p{\linewidth}}}]{
                \texttt{\tiny{[GM$|$GM]}}
                \texttt{SPRICH: Ich habe keine Aktion ausgeführt, da meine Objekte bereits an den Zielpositionen sind. Bitte teile mir mit, ob dein I erfolgreich nach (4,7) bewegt wurde.} \\
            }
        }
    }
    & & \\ \\

    \theutterance \stepcounter{utterance}  
    & & \multicolumn{4}{p{0.6\linewidth}}{
        \cellcolor[rgb]{0.9,0.9,0.9}{
            \makecell[{{p{\linewidth}}}]{
                \texttt{\tiny{[P1$\langle$GM]}}
                \texttt{'I' erfolgreich nach (4,7) verschoben. Dein Raster sieht jetzt so aus:} \\
\texttt{```} \\
\texttt{ 1234567} \\
\texttt{╔═╤═╤═══╗} \\
\texttt{║H│◌│◌◌◌║ 1} \\
\texttt{║◌├─┤◌◌◌║ 2} \\
\texttt{╟─┤◌├─┬─╢ 3} \\
\texttt{║◌│◌│◌│◌║ 4} \\
\texttt{╟─┼─┘◌│◌║ 5} \\
\texttt{║◌│C◌┌┴─╢ 6} \\
\texttt{║T│WI│◌◌║ 7} \\
\texttt{╚═╧══╧══╝} \\
\texttt{```} \\
\texttt{Ihr seid gerade bei Runde 12 von maximal 20.} \\
\texttt{Ihr habt gemeinsam 5 von maximal 12 Strafpunkten angesammelt.} \\
\texttt{Dein Mitspieler hat dir diese Nachricht geschickt:} \\
\\ 
\\ 
\texttt{"Ich habe keine Aktion ausgeführt, da meine Objekte bereits an den Zielpositionen sind. Bitte teile mir mit, ob dein I erfolgreich nach (4,7) bewegt wurde."} \\
\\ 
\\ 
\texttt{Wie lautet dein nächster Befehl?} \\
            }
        }
    }
    & \\ \\

    \theutterance \stepcounter{utterance}  
    & \multicolumn{4}{p{0.6\linewidth}}{
        \cellcolor[rgb]{0.9,0.9,0.9}{
            \makecell[{{p{\linewidth}}}]{
                \texttt{\tiny{[P1$\rangle$GM]}}
                \texttt{SPRICH: Ja, I wurde erfolgreich nach (4,7) bewegt. Alle Objekte sind nun an ihren Zielpositionen. Ich bin bereit, das Spiel zu beenden.} \\
            }
        }
    }
    & & \\ \\

    \theutterance \stepcounter{utterance}  
    & & & \multicolumn{2}{p{0.3\linewidth}}{
        \cellcolor[rgb]{0.9,0.9,0.9}{
            \makecell[{{p{\linewidth}}}]{
                \texttt{\tiny{[GM$|$GM]}}
                \texttt{SPRICH: Ja, I wurde erfolgreich nach (4,7) bewegt. Alle Objekte sind nun an ihren Zielpositionen. Ich bin bereit, das Spiel zu beenden.} \\
            }
        }
    }
    & & \\ \\

    \theutterance \stepcounter{utterance}  
    & & \multicolumn{4}{p{0.6\linewidth}}{
        \cellcolor[rgb]{0.9,0.9,0.9}{
            \makecell[{{p{\linewidth}}}]{
                \texttt{\tiny{[GM$\rangle$P2]}}
                \texttt{Deine Nachricht wurde an deinen Mitspieler übermittelt.} \\
\\ 
\\ 
\texttt{Ihr seid gerade bei Runde 13 von maximal 20.} \\
\texttt{Ihr habt gemeinsam 5 von maximal 12 Strafpunkten angesammelt.} \\
\texttt{Dein Mitspieler hat dir diese Nachricht geschickt:} \\
\\ 
\\ 
\texttt{"Ja, I wurde erfolgreich nach (4,7) bewegt. Alle Objekte sind nun an ihren Zielpositionen. Ich bin bereit, das Spiel zu beenden."} \\
\\ 
\\ 
\texttt{Wie lautet dein nächster Befehl?} \\
            }
        }
    }
    & \\ \\

    \theutterance \stepcounter{utterance}  
    & & & \multicolumn{4}{p{0.6\linewidth}}{
        \cellcolor[rgb]{0.9,0.9,0.9}{
            \makecell[{{p{\linewidth}}}]{
                \texttt{\tiny{[GM$\langle$P2]}}
                \texttt{SPRICH: fertig!} \\
            }
        }
    }
     \\ \\

    \theutterance \stepcounter{utterance}  
    & & & \multicolumn{2}{p{0.3\linewidth}}{
        \cellcolor[rgb]{0.9,0.9,0.9}{
            \makecell[{{p{\linewidth}}}]{
                \texttt{\tiny{[GM$|$GM]}}
                \texttt{SPRICH: fertig!} \\
            }
        }
    }
    & & \\ \\

    \theutterance \stepcounter{utterance}  
    & & \multicolumn{4}{p{0.6\linewidth}}{
        \cellcolor[rgb]{0.9,0.9,0.9}{
            \makecell[{{p{\linewidth}}}]{
                \texttt{\tiny{[P1$\langle$GM]}}
                \texttt{Deine Nachricht wurde an deinen Mitspieler übermittelt.} \\
\\ 
\\ 
\texttt{Ihr seid gerade bei Runde 13 von maximal 20.} \\
\texttt{Ihr habt gemeinsam 5 von maximal 12 Strafpunkten angesammelt.} \\
\texttt{Dein Mitspieler hat dir diese Nachricht geschickt:} \\
\\ 
\\ 
\texttt{"fertig!"} \\
\\ 
\\ 
\texttt{Wie lautet dein nächster Befehl?} \\
            }
        }
    }
    & \\ \\

    \theutterance \stepcounter{utterance}  
    & \multicolumn{4}{p{0.6\linewidth}}{
        \cellcolor[rgb]{0.9,0.9,0.9}{
            \makecell[{{p{\linewidth}}}]{
                \texttt{\tiny{[P1$\rangle$GM]}}
                \texttt{SPRICH: fertig!} \\
            }
        }
    }
    & & \\ \\

    \theutterance \stepcounter{utterance}  
    & & & \multicolumn{2}{p{0.3\linewidth}}{
        \cellcolor[rgb]{0.9,0.9,0.9}{
            \makecell[{{p{\linewidth}}}]{
                \texttt{\tiny{[GM$|$GM]}}
                \texttt{SPRICH: fertig!} \\
            }
        }
    }
    & & \\ \\

    \theutterance \stepcounter{utterance}  
    & & & \multicolumn{2}{p{0.3\linewidth}}{
        \cellcolor[rgb]{0.9,0.9,0.9}{
            \makecell[{{p{\linewidth}}}]{
                \texttt{\tiny{[GM$|$GM]}}
                \texttt{true} \\
            }
        }
    }
    & & \\ \\

    \theutterance \stepcounter{utterance}  
    & & & \multicolumn{2}{p{0.3\linewidth}}{
        \cellcolor[rgb]{0.9,0.9,0.9}{
            \makecell[{{p{\linewidth}}}]{
                \texttt{\tiny{[GM$|$GM]}}
                \texttt{Initial states:} \\
\texttt{```} \\
\texttt{Player 1:} \\
\texttt{ 1234567} \\
\texttt{╔═╤═╤═══╗} \\
\texttt{║C│◌│◌TI║ 1} \\
\texttt{║◌├─┤W◌◌║ 2} \\
\texttt{╟─┤◌├─┬─╢ 3} \\
\texttt{║◌│◌│◌│◌║ 4} \\
\texttt{╟─┼─┘◌│◌║ 5} \\
\texttt{║◌│◌H┌┴─╢ 6} \\
\texttt{║◌│◌◌│◌◌║ 7} \\
\texttt{╚═╧══╧══╝} \\
\\ 
\\ 
\texttt{Player 2:} \\
\texttt{ 1234567} \\
\texttt{╔═╤═╤═══╗} \\
\texttt{║H│◌│◌◌◌║ 1} \\
\texttt{║◌├─┤◌◌◌║ 2} \\
\texttt{╟─┤◌├─┬─╢ 3} \\
\texttt{║◌│◌│C│◌║ 4} \\
\texttt{╟─┼─┘◌│◌║ 5} \\
\texttt{║◌│◌◌┌┴─╢ 6} \\
\texttt{║T│WI│◌◌║ 7} \\
\texttt{╚═╧══╧══╝} \\
\\ 
\texttt{```} \\
            }
        }
    }
    & & \\ \\

    \theutterance \stepcounter{utterance}  
    & & & \multicolumn{2}{p{0.3\linewidth}}{
        \cellcolor[rgb]{0.9,0.9,0.9}{
            \makecell[{{p{\linewidth}}}]{
                \texttt{\tiny{[GM$|$GM]}}
                \texttt{End states:} \\
\texttt{```} \\
\texttt{Player 1:} \\
\texttt{ 1234567} \\
\texttt{╔═╤═╤═══╗} \\
\texttt{║H│◌│◌◌◌║ 1} \\
\texttt{║◌├─┤◌◌◌║ 2} \\
\texttt{╟─┤◌├─┬─╢ 3} \\
\texttt{║◌│◌│◌│◌║ 4} \\
\texttt{╟─┼─┘◌│◌║ 5} \\
\texttt{║◌│C◌┌┴─╢ 6} \\
\texttt{║T│WI│◌◌║ 7} \\
\texttt{╚═╧══╧══╝} \\
\\ 
\\ 
\texttt{Player 2:} \\
\texttt{ 1234567} \\
\texttt{╔═╤═╤═══╗} \\
\texttt{║H│◌│◌◌◌║ 1} \\
\texttt{║◌├─┤◌◌◌║ 2} \\
\texttt{╟─┤◌├─┬─╢ 3} \\
\texttt{║◌│◌│◌│◌║ 4} \\
\texttt{╟─┼─┘◌│◌║ 5} \\
\texttt{║◌│C◌┌┴─╢ 6} \\
\texttt{║T│WI│◌◌║ 7} \\
\texttt{╚═╧══╧══╝} \\
\\ 
\texttt{```} \\
            }
        }
    }
    & & \\ \\

    \theutterance \stepcounter{utterance}  
    & & & \multicolumn{2}{p{0.3\linewidth}}{
        \cellcolor[rgb]{0.9,0.9,0.9}{
            \makecell[{{p{\linewidth}}}]{
                \texttt{\tiny{[GM$|$GM]}}
                \texttt{* success: True} \\
\texttt{* lose: False} \\
\texttt{* aborted: False} \\
\texttt{{-}{-}{-}{-}{-}{-}{-}} \\
\texttt{* Shifts: 4.00} \\
\texttt{* Max Shifts: 8.00} \\
\texttt{* Min Shifts: 4.00} \\
\texttt{* End Distance Sum: 0.00} \\
\texttt{* Init Distance Sum: 30.73} \\
\texttt{* Expected Distance Sum: 20.95} \\
\texttt{* Penalties: 5.00} \\
\texttt{* Max Penalties: 12.00} \\
\texttt{* Rounds: 13.00} \\
\texttt{* Max Rounds: 20.00} \\
\texttt{* Object Count: 5.00} \\
            }
        }
    }
    & & \\ \\

\end{supertabular}
}

\end{document}
